% concatenated from design___m-estimation.tex
%\input{tcilatex}
%\setcounter{MaxMatrixCols}{10}


\documentclass[times, authoryear, 12pt, twoside]{elsarticle}
%%%%%%%%%%%%%%%%%%%%%%%%%%%%%%%%%%%%%%%%%%%%%%%%%%%%%%%%%%%%%%%%%%%%%%%%%%%%%%%%%%%%%%%%%%%%%%%%%%%%%%%%%%%%%%%%%%%%%%%%%%%%%%%%%%%%%%%%%%%%%%%%%%%%%%%%%%%%%%%%%%%%%%%%%%%%%%%%%%%%%%%%%%%%%%%%%%%%%%%%%%%%%%%%%%%%%%%%%%%%%%%%%%%%%%%%%%%%%%%%%%%%%%%%%%%%
\usepackage{geometry}
\usepackage[labelfont=bf]{caption}
\usepackage{amsmath,amsfonts,amssymb,amsthm,color,epsfig,graphicx,url, eurosym, latexsym}
\usepackage{appendix}
\usepackage[utf8]{inputenc}
\usepackage[english]{babel}

\setcounter{MaxMatrixCols}{10}
%TCIDATA{OutputFilter=LATEX.DLL}
%TCIDATA{Version=5.00.0.2606}
%TCIDATA{<META NAME="SaveForMode" CONTENT="1">}
%TCIDATA{BibliographyScheme=BibTeX}
%TCIDATA{LastRevised=Wednesday, April 22, 2026 18:12:58}
%TCIDATA{<META NAME="GraphicsSave" CONTENT="32">}
%TCIDATA{Language=American English}

\urldef\myurl\url{www.ualberta.ca/~dwiens/}
\pagestyle{myheadings}
\markboth{Minimax Robust Designs for M-Estimated Models}{Rui Hu, Douglas P. Wiens}
\textheight=22cm
\newtheorem{theorem}{Theorem}
\newtheorem{corollary}{Corollary}
\newtheorem{lemma}{Lemma}
\newtheorem{remark}{Remark}
\newtheorem{Discussion}{Discussion}
\makeatletter
\def\ps@pprintTitle{  \let\@oddhead\@empty
  \let\@evenhead\@empty
  \def\@oddfoot{\reset@font\hfil\thepage\hfil}
  \let\@evenfoot\@oddfoot
}


\begin{document}

\def\T{{ \mathrm{\scriptscriptstyle T} }}

\date{\today }

\input{amssym.def}

\begin{frontmatter}
\title{Minimax Robust Designs for M-Estimated Models}

\author[A1]{Rui Hu\corref{mycorrespondingauthor1}}
\address[A1]{Mathematics \& Statistics,
	MacEwan University,
	Edmonton, Canada,  T5J 4S2}

\author[A2]{Douglas P. Wiens\corref{mycorrespondingauthor2}}
\address[A2]{Mathematical \& Statistical Sciences,
	University of Alberta,
	Edmonton, Canada,  T6G 2G1}

\cortext[mycorrespondingauthor1]{E-mail: \url{hur3@macewan.ca}. } 
\cortext[mycorrespondingauthor2]{E-mail: \url{doug.wiens@ualberta.ca}. Computing code is at {\myurl}.} 

\begin{abstract}
Experimental designs that are minimax in the presence of
model misspecifications have been constructed so as to minimize the maximum,
over classes of alternate response models, of the integrated mean squared
error of the predicted values. The theory to date has focussed almost exclusively on Least Squares estimates. Here we extend this theory to designs tailored for M-estimation of parameters, thus obtaining additional robustness against outlying, or otherwise erroneous, responses. We show that, subject to a minor change in a tuning constant, designs optimal for Least Squares remain so asymptotically for M-estimation. We argue that even this minor change should be ignored, and the tuning constant chosen in an \textit{ad hoc} but sensible manner that does not depend on which M-estimate is being employed. A rather surprising additional result is that our designs and estimates, derived under an assumption of i.i.d. errors, are also robust, in a minimax sense, against broad classes of correlation structures.
\end{abstract}

\begin{keyword} %alphabetical order
asymptotics \sep
finite design space \sep
misspecified model \sep
regression design \sep
robustness against dependence. 
%\MSC[2010] Primary 62F35 \sep Secondary 62K05
\end{keyword}
\end{frontmatter}

\section{Introduction and summary \label{section: intro}}

The theory and practice of robustness of design, for possibly misspecified
response functions, is well-developed as it applies to cases in which
parameter estimation is to be carried out by Least Squares (\textsc{ls}). See \cite{w15} for a review. An investigator seeking model robustness might naturally be concerned as well
with robustness against outlying data points, or more generally against a
misspecified data-generating probability distribution, and hence seek
M-estimates of the parameters. Indeed, as pointed out by \cite{m97}, `a robust inference will profit from an optimal design and ... an optimal design is more reasonable if it allows also a robust analysis of the data'. See also \cite{mm01} for the special case of $L_1$ estimation.

\cite{w94,w96a}
studied this design problem for quite limited classes of approximate
responses on continuous design spaces and obtained asymptotic results under
rather restrictive conditions. \cite{ww10} carried out a small simulation
study and found that there was little apparent dependence of the designs on the
method of estimation.

In this article we strengthen and extend these findings. We show that, if the
parameters of the assumed model are to be estimated by ordinary
M-estimation, then designs optimal, in the sense of minimizing the maximum
mean squared error of the predictions, for \textsc{ls} remain so,
asymptotically, for M-estimation.

We note that M-estimates offer protection against outlying responses, but
not necessarily against outliers in the factor space. This latter type of
protection is furnished by Generalized M-estimation. But since our focus
centres on design points chosen by the experimenter, such outlyingness is
not an issue.

In \S \ref{section: asymptotics} of this article we present the asymptotic
theory on which our design problem will be based. The challenges there are
somewhat unique, since we do not assume that the fitted response model is the correct
one, and must allow for a broad class of alternatives. The ensuing Theorem \ref{thm:main} seems to be the most explicit and comprehensive response to this challenge, in the literature. 

In \S \ref%
{section: minimax theory} we address the design problem. We show that the
optimally robust designs depend on the anticipated method of estimation only
through a tuning constant. Then in \S \ref{section: dependence} we present a surprising (to us) 
result under which the designs and estimates are also minimax robust against
broad classes of correlation structures; a specific case is that in which
the random errors are equicorrelated.

Examples and methods of implementation are studied and discussed in \S \ref%
{section: examples}. We argue there that, although the aforementioned tuning
constant depends on the method of estimation through unknown parameters,
this dependence is so slight that it should be ignored and the tuning
constant chosen in an \textit{ad hoc} but sensible manner -- as the relative
emphasis placed by the designer on bias reduction versus variance reduction.
The result is then that the designs are completely independent of which
M-estimate is to be used.

Proofs are in the Appendix. The \textsc{matlab} code used to prepare the
examples is available on the second author's personal website.

\section{Minimax robustness of design \label{section: minimax}}

\subsection{Asymptotic theory \label{section: asymptotics}}

Our minimax design problem is phrased in terms of an approximate regression
response 
\begin{equation}
E\left[ Y\left( \boldsymbol{x}\right) \right] \approx \boldsymbol{f}^{\T
}\left( \boldsymbol{x}\right) \boldsymbol{\theta },  \label{approx}
\end{equation}%
for $p$ regressors $\boldsymbol{f}$, each functions of vectors $\boldsymbol{x%
}$ of $q$ independent variables, ranging over a finite \textit{design space} 
$\chi=\left\{ \boldsymbol{x}_{1},...,\boldsymbol{x}_{N}\right\}
\subset \mathbb{R}^{q}$, and for a parameter vector $\boldsymbol{\theta }%
_{p\times 1}$. 

Since (\ref{approx}) is an approximation, the interpretation of $\boldsymbol{%
\theta }$ is unclear; we \textit{define }this target parameter by 
\begin{equation}
\boldsymbol{\theta }_{0}=\arg \min_{\boldsymbol{\eta }}\sum_{i=1}^{N}\left( E%
\left[ Y\left( \boldsymbol{x}_{i}\right) \right] -\boldsymbol{f}^{\T
}\left( \boldsymbol{x}_{i}\right) \boldsymbol{\eta }\right) ^{2}.
\label{theta def}
\end{equation}%
Equivalently, and with model error $\tau \left( \boldsymbol{x}\right) \overset{def}{=}E%
\left[ Y\left( \boldsymbol{x}\right) \right] -\boldsymbol{f}^{\T }\left( 
\boldsymbol{x}\right) \boldsymbol{\theta }_{0}$, we have 
\begin{equation}
\sum_{i=1}^{N}\boldsymbol{f}\left( \boldsymbol{x}_{i}\right) \tau \left( 
\boldsymbol{x}_{i}\right) =\boldsymbol{0}_p.\   \label{orthogonality}
\end{equation}%
Assuming that $\mathcal{X}$ is rich enough that the matrix\ $\boldsymbol{A}%
\overset{def}{=}\sum_{i=1}^{N}\boldsymbol{f}\left( \boldsymbol{x}_{i}\right) 
\boldsymbol{f}^{\T }\left( \boldsymbol{x}_{i}\right) $ is invertible,
the parameter defined by (\ref{theta def}) and (\ref{orthogonality}) is
unique. In order that variance and bias remain of a comparable order
asymptotically, we bound the approximation error in (\ref{approx}) by
assuming that%
\begin{equation}
\sum_{i=1}^{N}\tau ^{2}\left( \boldsymbol{x}_{i}\right) \leq \kappa ^{2}/n,
\label{bound}
\end{equation}%
for a constant $\kappa $. Here $n$ is the total number of experimental runs to be carried out.

Our model is thus given by $E\left[ Y\left( \boldsymbol{x}\right) \right] =%
\boldsymbol{f}^{\T }\left( \boldsymbol{x}\right) \boldsymbol{\theta }%
_{0}+\tau \left( \boldsymbol{x}\right) $, for an unknown model error $\tau
\left( \boldsymbol{\cdot }\right) $ constrained by (\ref{orthogonality}) and
(\ref{bound}). Let $\Upsilon $ be the class of such functions $\tau \left( 
\boldsymbol{\cdot }\right) $.

At such values of $\boldsymbol{x}$, observations $Y\left( \boldsymbol{x}%
\right) $ may be made, with additive random error: $Y\left( \boldsymbol{x}%
\right) =E\left[ Y\left( \boldsymbol{x}\right) \right] +\varepsilon ,$ for
i.i.d., symmetrically distributed errors $\varepsilon $ with finite variance. We make no further distributional assumptions. In particular, the distribution of $\varepsilon $ may be different from that assumed by the experimenter, thus calling for the type of robustness furnished by M-estimation.

Thus, given observations $\left\{ Y_{j}\left( \boldsymbol{x}_{i}\right)
|j=1,...,n_{i}\right\} $, with $n_{i}$ replicates at distinct points $\boldsymbol{x}_{i}$, we suppose
that $\boldsymbol{\theta }_{0}$ will be estimated by M-estimation with an
auxiliary estimate of scale. For a function $\psi $ with properties as in
C1) below, the estimate satisfies 
\begin{equation}
\sum_{i,j}\psi \left( \frac{Y_{j}\left( \boldsymbol{x}%
_{i}\right) -\boldsymbol{f}^{\T }\left( \boldsymbol{x}_{i}\right) 
\boldsymbol{\hat{\theta}}_{n}}{\hat{\sigma}_{n}}\right) \boldsymbol{f}\left( 
\boldsymbol{x}_{i}\right)=\boldsymbol{0}_p \label{regression},
\end{equation}%
where $\hat{\sigma}_{n}$ is a consistent estimate of scale, for instance 
\begin{equation}
\hat{\sigma}_{n}=\left. \text{\textsc{median}}\left\{ \left\vert Y_{j}\left( 
\boldsymbol{x}_{i}\right) -\boldsymbol{f}^{\T }\left( \boldsymbol{x}%
_{i}\right) \boldsymbol{\hat{\theta}}_{n}\right\vert \right\} \right/ \Phi
^{-1}\left( .75\right),  \label{scale}
\end{equation}%
which is consistent for the standard deviation $\sigma $ if the data are
Normal. Under the mild conditions of \cite{rs86}, $\boldsymbol{\hat{\theta}}%
_{n}$ has the same asymptotic properties as if $\hat{\sigma}_{n}$ were
replaced by $\sigma $. Thus, with $\psi _{\sigma }\left( x\right) \overset{%
def}{=}\psi \left( x/\sigma \right) $, we have that
\begin{equation}
\sum_{i,j}\psi _{\sigma }\left( Y_{j}\left( \boldsymbol{x}%
_{i}\right) -\boldsymbol{f}^{\T }\left( \boldsymbol{x}_{i}\right) 
\boldsymbol{\hat{\theta}}_{n}\right) \boldsymbol{f}\left( \boldsymbol{x}%
_{i}\right)=\boldsymbol{0}_p,  \label{m-eqn}
\end{equation}%
and we replace $\sigma $ by $\hat{\sigma}_{n}$ in the applications. For
the computations one would iterate between (\ref{regression}) and (\ref%
{scale}). See \cite{h81} for further details of M-estimation, \cite{hw17} for applications to testing and \cite{h18} for applications in spatial sampling.

There is a rich literature concerning minimax variance M-estimation when the true error distribution is assumed to vary over a neighbourhood of the assumed distribution -- see for instance \cite{h64}, \cite{w86}, \cite{cw85}, \cite{cw89} for examples of solutions in various such neighbourhoods. We do not pursue this further, and leave the choice of a specific M-estimator up to the experimenter.

By a \textit{design} we mean a probability mass function $\xi$ on $\mathcal{X}$, to be approximated as necessary by an implementable, $n$-point design $\xi_{n}$, with $\xi_{n}(\boldsymbol{x}_{i})=n_{i}/n$. For any design $\xi$ define 
\begin{equation*}
\boldsymbol{M}_{0}\left( \xi \right) =\sum_{i=1}^{N}\boldsymbol{f}%
\left( \boldsymbol{x}_{i}\right) \boldsymbol{f}^{\T }\left( \boldsymbol{x}_{i}%
\right) \xi \left( \boldsymbol{x}_{i}\right) \text{ and }\boldsymbol{b}%
_{0}\left( \xi \right) =\sum_{i=1}^{N}\boldsymbol{f}%
\left( \boldsymbol{x}_{i}\right) \tau\left( \boldsymbol{x}_{i}%
\right) \xi \left( \boldsymbol{x}_{i}\right) . 
\end{equation*}

We define $\sigma _{M}^{2}=\left. E\left[ \psi _{\sigma }^{2}\left(
\varepsilon \right) \right] \right/ \left( E\left[ \psi _{\sigma }^{\prime
}\left( \varepsilon \right) \right] \right) ^{2}$ (assumed finite), 
and for a design $\xi _{n}$ we set 
$\boldsymbol{M}_{n}=\boldsymbol{M}_{0}\left( \xi _{n}\right) $, and 
$\boldsymbol{b}_{n}=\boldsymbol{b}_{0}\left( \xi _{n}\right) $. In the
Appendix we prove Theorem \ref{thm:main}, stated at the end of this section,
which makes precise the asymptotic normality of $\boldsymbol{\hat{\theta}}%
_{n}-\boldsymbol{\theta }_{0}$:%
\begin{equation}
\boldsymbol{\hat{\theta}}_{n}-\boldsymbol{\theta }_{0}\sim AN\left( 
\boldsymbol{M}_{n}^{-1}\boldsymbol{b}_{n},\left( \sigma
_{M}^{2}/n\right) \boldsymbol{M}_{n}^{-1}\right) .  \label{AN}
\end{equation}%
\noindent For this we define $\boldsymbol{\tilde{\theta}}_{n}=\boldsymbol{%
\theta }_{0}+\boldsymbol{M}_{n}^{-1}\boldsymbol{b}_{n}$, and will show
that $\sqrt{n}\left( \boldsymbol{\hat{\theta}}_{n}-\boldsymbol{\theta }_{0}-%
\boldsymbol{M}_{n}^{-1}\boldsymbol{b}_{n}\right) =\sqrt{n}\left( 
\boldsymbol{\hat{\theta}}_{n}-\boldsymbol{\tilde{\theta}}_{n}\right) $ is
asymptotically normal with mean $\boldsymbol{0}_p$.

For the proof of Theorem \ref{thm:main} we make the following assumptions
and definitions. We denote the smallest and largest eigenvalues of a matrix
by $ch_{\min }$ and $ch_{\max }$ respectively.

\begin{description}
\item[C1)] The function $\psi \left( \cdot \right) $ is weakly increasing,
twice differentiable and odd. We define $m_{1}=\sup_t \psi _{\sigma }^{\prime }\left(t\right) $ and $m_{2}=\sup_t|\psi _{\sigma }^{\prime \prime }\left(t\right) |$, both assumed finite.

\item[C2)] The sequence $A_{n}=ch_{\min }\boldsymbol{M}_{n}$ is bounded
away from zero.

\item[C3)] The sequence of designs has a weak limit $\xi _{\ast }$: $\xi _{n}%
\overset{d}{\rightarrow }\xi _{\ast }$ as $n\rightarrow \infty $.

\begin{remark}
\label{remark: assumptions}
Assumption C1) is standard, but precludes the use of redescenders (\cite{m04}). It implies that $E\left[ \psi \left( \varepsilon \right) %
\right] =E\left[ \psi ^{\prime \prime }\left( \varepsilon \right) \right] =0$%
, since the distribution of $\varepsilon $ is symmetric, and that $E\left[ \psi
^{\prime }\left( \varepsilon \right) \right] >0$, else $\psi \left( \cdot
\right) $ is identically zero. For C2), that $ch_{\min }\boldsymbol{M}_{n} >>0$ ensures that $\boldsymbol{M}_{n}$ remains invertible,
and is a natural property of the design. For C3), that the designs have a
weak limit asserts only that the design weights $\xi_{n}(\boldsymbol{x}_{i})$ converge, which is
a requirement of our numerical algorithm, discussed in \S \ref{section:
construction}. The limit will be used in the proof of Theorem \ref{thm:main}
when we invoke Theorem 5.9 of \cite{vw98}, upon which we base our proof that 
$\boldsymbol{\hat{\theta}}_{n}-\boldsymbol{\tilde{\theta}}_{n}$ is a $\sqrt{n%
}$-consistent estimate of zero. 
\end{remark}
\end{description}

\begin{theorem}
\label{thm:main}With notation as above, we have that, as $n\rightarrow
\infty $, 
\begin{equation*}
\sqrt{n}\boldsymbol{M}_{n}^{1/2}\left( \boldsymbol{\hat{\theta}}_{n}-%
\boldsymbol{\theta }_{0}-\boldsymbol{M}_{n}^{-1}\boldsymbol{b}%
_{n}\right) \overset{d}{\longrightarrow }N\left( \boldsymbol{0}_p,\sigma
_{M}^{2}\boldsymbol{I}_{p}\right) .
\end{equation*}
\end{theorem}

\subsection{Minimax design theory \label{section: minimax theory}}

We define our loss in terms of the Integrated Mean Squared Error of the
predictors $\hat{Y}\left( \boldsymbol{x}\right) =\boldsymbol{f}^{\T
}\left( \boldsymbol{x}\right) \boldsymbol{\hat{\theta}}_{n}$: 
\begin{equation}
\text{\textsc{imse}}\left( \xi |\tau \right) =\sum_{\boldsymbol{x}\mathbb{
\in }\chi}E\left\{ \left( E\left[ Y\left( \boldsymbol{x}\right) %
\right] -\hat{Y}\left( \boldsymbol{x}\right) \right) ^{2}\right\} =\sum_{%
\boldsymbol{x}\mathbb{\in }\chi}E\left\{ \left( \tau \left( 
\boldsymbol{x}\right) -\boldsymbol{f}^{\T }\left( \boldsymbol{x}\right)
\left( \boldsymbol{\hat{\theta}}_{n}\left( \xi |\tau \right) -\boldsymbol{\theta }_{0}(\tau)%
\right) \right) ^{2}\right\} .  \label{IM}
\end{equation}%
We will maximize an asymptotic approximation to \textsc{imse}$\left( \xi |\tau \right) $ over $\tau $
satisfying (\ref{orthogonality}) and (\ref{bound}), and then find designs
minimizing this maximum. The approximation uses the asymptotic mean and covariance given in (\ref{AN}). See \cite{l99} for a discussion of limiting versus asymptotic moments. 

Expanding (\ref{IM}) and using (\ref{orthogonality}) results
in 
\begin{equation}
\text{\textsc{imse}}\left( \xi|\tau \right)=\sum_{i=1}^{N}\tau ^{2}\left( \boldsymbol{x}%
_{i}\right) +\sum_{i=1}^{N}\boldsymbol{f}^{\T }\left( \boldsymbol{x}%
_{i}\right) E\left[ \left( \boldsymbol{\hat{\theta}}_{n}(\xi|\tau)-\boldsymbol{\theta }%
_{0}(\tau)\right)\left( \boldsymbol{\hat{\theta}}_{n}\left( \xi |\tau \right) -\boldsymbol{\theta }_{0}(\tau)%
%
\right) ^{\T }\right] \boldsymbol{f}\left( \boldsymbol{x}_{i}\right).
\notag 
\end{equation}
Substituting the asymptotic moments gives the asymptotic approximation
\begin{eqnarray}
\text{\textsc{aimse}}\left( \xi_n |\tau \right)&=&\sum_{i=1}^{N}\tau ^{2}\left( \boldsymbol{x}_{i}\right) +\sum_{i=1}^{N}%
\boldsymbol{f}^{\T }\left( \boldsymbol{x}_{i}\right) \left[ \boldsymbol{M%
}_{n}^{-1}\boldsymbol{b}_{n}\boldsymbol{b}_{n}^{\T }\boldsymbol{M}%
_{n}^{-1}+\left( \sigma _{M}^{2}/n\right) \boldsymbol{M}_{n}^{-1}\right] 
\boldsymbol{f}\left( \boldsymbol{x}_{i}\right) \notag
\\
&=&\sum_{i=1}^{N}\tau ^{2}\left( \boldsymbol{x}_{i}\right) +\boldsymbol{b}%
_{n}^{\T }\boldsymbol{M}_{n}^{-1}\mathbf{A}\boldsymbol{M}_{n}^{-1}%
\boldsymbol{b}_{n}+\left( \sigma _{M}^{2}/n\right) tr\boldsymbol{AM}%
_{n}^{-1} \text{.}  \label{IMSEn}
\end{eqnarray}

For the maximization it is convenient to introduce an
orthogonal basis for the space of regressors. Define $\boldsymbol{F}%
_{N\times p}=\left( \boldsymbol{f}\left( \boldsymbol{x}_{1}\right) ,\cdot
\cdot \cdot ,\boldsymbol{f}\left( \boldsymbol{x}_{N}\right) \right) ^{\T
}$, and for a design $\xi $ on $\chi $, $\boldsymbol{D}\left( \xi \right)
=diag\left( \xi _{1},\cdot \cdot \cdot ,\xi _{N}\right) $. By the
Gram-Schmidt process we can construct a matrix $\boldsymbol{Q}_{N\times p}$
whose orthonormal columns form a basis for the column space of $\boldsymbol{F%
}$ -- assumed to be of dimension $p$.

In the Appendix we prove the following theorem.

\begin{theorem}
\label{thm: max loss} For any design for which $\boldsymbol{M}_{0}\left( \xi \right)$ in non-singular, define $p\times p$ matrices 
\begin{equation*}
\boldsymbol{R}\left( \xi \right) =\boldsymbol{Q}^{\T }\boldsymbol{D}%
\left( \xi \right) \boldsymbol{Q},\text{ }\boldsymbol{S}\left( \xi \right) =%
\boldsymbol{Q}^{\T }\boldsymbol{D}^{2}\left( \xi \right) \boldsymbol{Q},%
\text{ }\boldsymbol{U}\left( \xi \right) =\boldsymbol{R}^{-1}\left( \xi
\right) \boldsymbol{S}\left( \xi \right) \boldsymbol{R}^{-1}\left( \xi
\right) .
\end{equation*}%
Then $\max_{\tau \in \Upsilon } \textsc{aimse} \left( \xi _{n}|\tau \right) $
is given by $n^{-1}$ times 
\begin{equation}
J\left( \xi _{n}\right) =\sigma _{M}^{2}tr\boldsymbol{R}^{-1}\left( \xi
_{n}\right) +\kappa ^{2}ch_{\max }\boldsymbol{U}\left( \xi _{n}\right) .
\label{I3}
\end{equation}
\end{theorem}

A further maximization of the \textsc{aimse} is discussed in the next
section, after which, in \S \ref{section: construction}, we minimize this
maximum, thus obtaining the minimax designs.

\subsection{Robustness against dependence \label{section: dependence}}

It is brought out in the proof of Theorem \ref{thm:main} -- see (\ref%
{represent}) of the Appendix -- that the estimate is representable as 
\begin{equation*}
\boldsymbol{\Phi }_{n}\overset{def}{=}\sqrt{n}\boldsymbol{M}%
_{n}^{1/2}\left( \boldsymbol{\hat{\theta}}_{n}-\boldsymbol{\theta }_{0}-%
\boldsymbol{M}_{n}^{-1}\boldsymbol{b}_{n}\right) =\boldsymbol{K}%
_{p\times n}\boldsymbol{u}+\boldsymbol{k}_{p\times 1}+o_{p}\left( 1\right) ,
\end{equation*}%
where $\boldsymbol{K}$ and $\boldsymbol{k}$ are non-random and $\boldsymbol{u}=\left(
U_{1},...,U_{n}\right) ^{\T }$ with $U_{i}=\psi _{\sigma }\left(
\varepsilon _{i}\right) \left/ E\left[ \psi _{\sigma }^{\prime }\left(
\varepsilon \right) \right] \right. $. We have up to now assumed that the
errors $\varepsilon _{i}$ are i.i.d., so that as well the $U_{i}$ are
i.i.d., with \textsc{cov}$\left[ \boldsymbol{u}\right] =\sigma _{M}^{2}%
\boldsymbol{I}_{n}$ and, asymptotically, \textsc{cov}$\left[ \boldsymbol{\Phi }_{n}\right]
=\sigma _{M}^{2}\boldsymbol{KK}^{\T }$. We now investigate the effect on
the estimate if the $U_{i}$ are correlated, or are heteroscedastic.

Suppose then that \textsc{cov}$\left[ \boldsymbol{u}\right] =\boldsymbol{C}%
_{n\times n}$, so that the asymptotic covariance
structure of $\boldsymbol{\Phi }_{n}$ becomes \textsc{cov}$\left[ 
\boldsymbol{\Phi }_{n}|\boldsymbol{C}\right] =\boldsymbol{KCK}^{\T }$.
Note that \textsc{cov}$\left[ \boldsymbol{\Phi }_{n}|\boldsymbol{C}\right] $
is non-decreasing in the Loewner ordering: $\boldsymbol{C}_{1}\preceq 
\boldsymbol{C}_{2}\Rightarrow $ \textsc{cov}$\left[ \boldsymbol{\Phi }_{n}|%
\boldsymbol{C}_{1}\right] \preceq $ \textsc{cov}$\left[ \boldsymbol{\Phi }%
_{n}|\boldsymbol{C}_{2}\right] $. If loss is measured by $\mathcal{L}\left( 
\boldsymbol{C}\right) =\phi \left( \text{\textsc{cov}}\left[ \boldsymbol{%
\Phi }_{n}|\boldsymbol{C}\right] \right) $, where $\phi \left( \cdot \right) 
$ is a function, such as the trace, determinant, maximum eigenvalue etc.
that is itself non-decreasing in the Loewner ordering, then $\mathcal{L}%
\left( \boldsymbol{C}\right) $ is non-decreasing in this ordering. In the
particular problem at hand, $\phi \left( \Sigma \right) =\sum_{i=1}^{N}%
\boldsymbol{f}^{\T }\left( \boldsymbol{x}_{i}\right) \Sigma \boldsymbol{f%
}\left( \boldsymbol{x}_{i}\right) $ is the integrated variance of the predictors when evaluated at $\boldsymbol{\Sigma }=$   \textsc{cov}$\left[ \Phi _{n}| \boldsymbol{C}\right] $.

Suppose now that $\left\Vert \cdot \right\Vert _{M}$ is a matrix norm,
induced by a vector norm $\left\Vert \cdot \right\Vert _{V}$, i.e. $%
\left\Vert \boldsymbol{C}\right\Vert _{M}=\sup_{\left\Vert x\right\Vert _{V}%
\text{ }=1}\left\Vert \boldsymbol{Cx}\right\Vert _{V}.$ Special cases are
the spectral radius $\left\Vert \boldsymbol{C}\right\Vert _{E}$ -- this is
the maximum eigenvalue since $\boldsymbol{C}$ is a covariance matrix -- and
the maximum absolute row sum $\left\Vert \boldsymbol{C}\right\Vert
_{1}=\max_{i}\sum_{j}\left\vert c_{ij}\right\vert $.

The following lemma is given in \cite{w25}; we repeat it here for
convenience.

\begin{lemma}
\label{lemma: dependence}For $\eta ^{2}>0$, covariance matrix $\boldsymbol{C}
$ and induced matrix norm $\left\Vert \boldsymbol{C}\right\Vert _{M}$, define 
\begin{equation*}
\mathcal{C}_{M}=\left\{ \boldsymbol{C}\left\vert {}\right. \boldsymbol{C}%
\succeq 0\text{ and }\left\Vert \boldsymbol{C}\right\Vert _{M}\leq \eta
^{2}\right\} .
\end{equation*}%
For the norm $\left\Vert \mathbf{\cdot }\right\Vert _{E}$ an equivalent
definition is $\mathcal{C}_{E}=\{\boldsymbol{C}\mid 0\preceq \boldsymbol{C}%
\preceq \eta ^{2}\boldsymbol{I}_{n}\}$. Then \noindent (i) in any such class 
$\mathcal{C}_{M}$, $\max_{\mathcal{C}_{M}}\mathcal{L}\left( \boldsymbol{C}%
\right) =\mathcal{L}\left( \eta ^{2}\boldsymbol{I}_{n}\right) $, and \noindent (ii) if $%
\mathcal{C}^{\prime }$ $\mathcal{\subseteq C}_{M}$ and $\eta ^{2}\boldsymbol{I}_{n}\in 
\mathcal{C}^{\prime }$, then $\max_{\mathcal{C}^{\prime }}\mathcal{L}\left( 
\boldsymbol{C}\right) =\mathcal{L}\left( \eta ^{2}\boldsymbol{I}_{n}\right) $%
.
\end{lemma}

A consequence of (i) of this lemma is that if one is carrying out a
statistical procedure with loss function $\mathcal{L}\left( \boldsymbol{C}%
\right) $, then a version of the procedure that minimizes $\mathcal{L}\left(
\eta ^{2}\boldsymbol{I}_{n}\right) $ is \textit{minimax} as $\boldsymbol{C}$
varies over $\mathcal{C}_{M}$. By (ii) this remains true for subsets of $%
\mathcal{C}_{M}$ that contain $\eta ^{2}\boldsymbol{I}_{n}$. To apply this
result we need only ensure that $\eta ^{2}$ is large enough that $\mathcal{C}%
_{M}$ contains the departures, from independence or homoscedasticity, that
are of interest. If so, the designs and estimates of this article -- i.e.
those derived under an assumption of i.i.d. errors -- enjoy the additional
optimality property of minimizing the maximum value of $\mathcal{L}\left( 
\boldsymbol{C}\right) $, as the covariance matrix $\boldsymbol{C}$ of $%
\boldsymbol{u}$ varies over $\mathcal{C}_{M}$.

Lemma \ref{lemma: dependence} is unsatisfactory in this M-estimation context
since it is based on the covariance structure of $\left\{ \psi _{\sigma
}\left( \varepsilon _{i}\right) \right\} $, rather than that of $\left\{
\varepsilon _{i}\right\} $. In general there is no tractable relationship
between the two. There are however important exceptions. Suppose that
under the dependence structure the $\left\{ \varepsilon _{i}\right\} $ are 
\textit{exchangeable}, hence \textit{equicorrelated}. Then the $\left\{ \psi
_{\sigma }\left( \varepsilon _{i}\right) \right\} $ also have this
structure: for some $\rho $ (typically \textit{not} the correlation among
the $\left\{ \varepsilon _{i}\right\} $), and some $\alpha ^{2}$, 
\begin{equation*}
\text{\textsc{cov}}\left[ \boldsymbol{u}\right] =\boldsymbol{C}=\alpha
^{2}\left( \left( 1-\rho \right) \boldsymbol{I}_{n}+\rho \boldsymbol{1}_{n}%
\boldsymbol{1}_{n}^{\T }\right) .
\end{equation*}%
We impose bounds $\left\vert \rho \right\vert \leq $ $\rho _{\max }<1$ and $%
\alpha ^{2}\leq \alpha _{\max }^{2}$. We shall work with the norm $%
\left\Vert \boldsymbol{C}\right\Vert _{1}=\alpha ^{2}\left( 1+\left(
n-1\right) \left\vert \rho \right\vert \right) $, this is also $\left\Vert 
\boldsymbol{C}\right\Vert _{E}$ if $\rho \geq 0$. Set $\eta ^{2}=\alpha
_{\max }^{2}\left( 1+\left( n-1\right) \rho _{\max }\right) .$ Then $%
\left\Vert \boldsymbol{C}\right\Vert _{1}\leq \eta ^{2}$ and Lemma \ref%
{lemma: dependence} applies: an assumption of i.i.d. errors with \textsc{var}%
$\left[ U_i) \right] =$ $\eta ^{2}$%
, which for convenience we write as $\eta ^{2}=\eta _{0}^{2}\sigma _{M}^{2}$%
, is minimax within a class containing all such equicorrelated error
structures. 

By a similar treatment one can show robustness against dependence structures under which $\varepsilon_{i}$ and $\varepsilon_{j}$ are independent for $|i-j| > q$, as is the case for $MA(q)$ errors. The case of independent but heteroscedastic errors is clearly
covered as well, by $\alpha _{\max }^{2}\geq \max_{i}$\textsc{var}$\left[
\psi _{\sigma }\left( \varepsilon _{i}\right) \right] $. 

\section{\protect\bigskip Construction of minimax designs\label{section:
construction}}

With 
\begin{equation}
\nu \overset{def}{=}\kappa ^{2}/\left( \eta _{0}^{2}\sigma _{M}^{2}+\kappa
^{2}\right) ,  \label{nu}
\end{equation}%
the results of \S \ref{section: minimax theory} easily accommodate those of 
\S \ref{section: dependence}: the maximum, over both $\tau \in \Upsilon $
and $\boldsymbol{C}\in \mathcal{C}_{1}$, of (\ref{IMSEn}) is given by $\left( \eta _{0}^{2}\sigma _{M}^{2}+\kappa
^{2}\right) /n$ (which does not depend on the design) times%
\begin{equation}
I_{\nu }\left( \xi _{n}\right) =\left( 1-\nu \right) tr\boldsymbol{R}%
^{-1}\left( \xi _{n}\right) +\nu ch_{\max }\boldsymbol{U}\left( \xi
_{n}\right) .  \label{Inu}
\end{equation}%
If robustness against
dependence is not of interest then set $\eta _{0}^{2}$ $=1$.

A \textit{minimax} design is a minimizer of $I_{\nu }\left( \xi _{n}\right) $%
. For fixed $\nu $, $I_{\nu }\left( \xi _{n}\right) $ is precisely the value
minimized, in \cite{w}, to obtain minimax designs for \textsc{ls} estimates,
thus justifying our statement in \S 1 that such designs remain minimax
optimal, asymptotically, for M-estimation.

The minimization of $I_{\nu }\left( \xi _{n}\right) $ is carried out
sequentially, as described in Theorem 5 of \cite{w}. Briefly, given a
current $k$-point design $\xi _{k}$, the loss resulting from the addition of
a design point at $\boldsymbol{x}_{i}$ is expanded as 
\begin{equation}
I_{\nu }\left( \xi _{k+1}^{(i)}\right) =I_{\nu }\left( \xi _{k}\right)
-t_{k,i}/k+O\left( k^{-2}\right) ,  \label{expansion}
\end{equation}%
and then $\boldsymbol{x}_{(i)}$, with $\left( i\right) =\arg \max_{i}t_{k,i}$%
, is added to the design. This is carried out to convergence, yielding a
design $\xi $ on $\mathcal{\chi }$ with intended allocations $n_{i}=n\xi
\left( \boldsymbol{x}_{i}\right) $. Typically most $%
\xi \left( \boldsymbol{x}_{i}\right) $ are zero, but otherwise the $n_{i}$
are not integers. To obtain implementable designs $\xi_{n}$ we first round up the $%
n_{i}$ to $\left\lceil n\xi \left( \boldsymbol{x}_{i}\right) \right\rceil $,
whose sum exceeds $n$. The excess is decreased stepwise, by removing points
whose value of $t_{k,i}$ in (\ref{expansion}) is a minimum. This method
typically results in only a very small increase in the minimized value of $%
I_{\nu }\left( \xi _{n}\right) $.

\section{Examples, implementations and discussion \label{section: examples}}

To obtain the minimax design only $\nu $, at (\ref{Inu}), need be specified
by the user. The simplest and most natural way to do this is to view $\nu $
as expressing the emphasis on the reduction of losses due to bias, rather
than to variation. Its choice is then up to the user; quite typically $\nu
=.5$ is chosen. In this method there is NO difference in the minimax designs
for different M-estimates.

Another method is suggested by the definition of $\nu $ at (\ref{nu}).
Although the parameters involved in this definition would not be known to
the designer, it is of interest to see how their values could affect the
resulting designs.

We begin by seeing how much $\nu $ can change from its value under \textsc{ls%
}. The relationship between $\nu $ using Least Squares and $\nu $ using the
M-estimate is, with $\gamma \overset{def}{=}\eta _{0}\sigma /\kappa $, that 
\begin{equation*}
\nu _{\text{\textsc{ls}}}=\left( \gamma ^{2}+1\right) ^{-1}\text{ and }\nu _{%
\text{\textsc{m}}}=\left( \frac{\gamma ^{2}\sigma _{M}^{2}}{\sigma ^{2}}%
+1\right) ^{-1}.
\end{equation*}%
We assess these assuming that $\varepsilon /\sigma \sim N(0,1)$ and that $%
\psi \left( x\right) =xI\left( \left\vert x\right\vert \leq c\right)
+sign(x)cI\left( \left\vert x\right\vert >c\right) $ (\cite{h64}).

\begin{lemma}
\label{lemma: maxdiff} With notation as above, 
\begin{equation*}
0\leq \nu _{\text{\textsc{ls}}}-\nu _{\text{\textsc{m}}}\leq \frac{\sqrt{\pi
/2}-1}{\sqrt{\pi /2}+1}\simeq .1124.
\end{equation*}%
The lower bound is attained only when the M-estimate is the \textsc{lse},
and the upper bound is attained when the M-estimate is the $L_{1}$ estimate
and $\gamma ^{2}=1/\sqrt{\pi /2}\simeq .7979$. At the maximum 
\begin{equation}
\nu _{\text{\textsc{ls}}}=\frac{\sqrt{\pi /2}}{\sqrt{\pi /2}+1}\simeq
.5562,\quad \nu _{\text{\textsc{m}}}=\frac{1}{\sqrt{\pi /2}+1}=1-\nu _{\text{%
\textsc{ls}}}\simeq .4438.  \label{nus}
\end{equation}
\end{lemma}

\begin{figure}[tbp]
%[tb]
\centering
\includegraphics[scale=1.2]{diff}
\caption{Differences $\protect\nu _{\text{\textsc{ls}}}-\protect\nu _{\text{%
\textsc{m}}}$ in terms of $\protect\gamma $ and $c$. }
\label{fig:diff}
\end{figure}

\begin{figure}[tbp]
%[tb]
\centering
\includegraphics[scale=1.2]{lindes}
\caption{Designs for linear regression; n=10, N=20. (a)-(c): Continuous
weights; (d)-(f): Integer allocations.}
\label{fig:lindes}
\end{figure}

\begin{figure}[tbp]
%[tb]
\centering
\includegraphics[scale=1.2]{cubdes}
\caption{Designs for cubic regression; n=20, N=20. (a)-(c): Continuous
weights; (d)-(f): Integer allocations.}
\label{fig:cubdes}
\end{figure}

See Figure \ref{fig:diff}. We have constructed designs, using the `worst
case' values $\nu _{\text{\textsc{ls}}}$ and $\nu _{\text{\textsc{m}}}$ of (%
\ref{nus}), and also their midpoint $\nu =.5$. See Figure \ref{fig:lindes}
for approximate linear regression and Figure \ref{fig:cubdes} for
approximate cubic regression, each on design spaces consisting of 20 equally
spaced points spanning $[-1,1]$. Given that $\nu _{\text{\textsc{ls}}}$, $%
\nu _{\text{\textsc{m}}}$ and the M-estimate were chosen so as to make the
designs as different as possible, these designs are remarkably similar --
even identical after being made implementable. The implementations all can
be described as taking the replicates that would otherwise be assigned by
the classically ($\nu =0$) I-optimal designs and spreading them out into
clusters at nearby design points. The I-optimal design for linear regression
places mass of $.5$ at each of $\pm 1$. That for cubic regression was
derived by \cite{s77} and places masses of $.1545$ and $.3455$ at $\pm 1$
and $\pm .4472$.

We conclude that an experimenter should feel quite safe in using the same
design for an experiment regardless of which M-estimate is to be employed,
and in choosing $\nu $ to represent his desired emphasis on bias reduction,
as posited at the beginning of this section. As well, the assumption of
i.i.d. errors is safe, even optimal in a minimax sense, amongst
equicorrelated or heteroscedastic error structures (and others).

\appendix

\section*{Appendix: Proofs}
%\setcounter{page}{1}
\setcounter{equation}{0} \renewcommand{\theequation}{A.\arabic{equation}} %
%\setcounter{lemma}{0} \renewcommand{\thelemma}{A.
%\arabic{lemma}} 
\setcounter{subsection}{0} \renewcommand{\thesubsection}{A.%
\arabic{subsection}}
%\setcounter{remark}{0}

\begin{remark}
It is convenient to identify a design with its design measure -- a discrete probability measure $\xi
_{n}\left( d\boldsymbol{x}\right) $ on $\mathcal{X}$. Thus, as an example,
if the design calls for $n_{i}$ runs to be made at $\boldsymbol{x}%
_{i}\in \mathcal{\chi }$, we write averages of functions $\phi $ of the data
as 
\begin{equation*}
\frac{1}{n}\sum_{\left\{ i|n_{i}>0\right\} }\sum_{j=1}^{n_{i}}\phi \left(
Y_{j}\left( \boldsymbol{x}_{i}\right) \right) =\sum_{\left\{
i|n_{i}>0\right\} }\frac{n_{i}}{n}\bar{\phi}\left( \boldsymbol{x}_{i}\right)
=\int_{\mathcal{X}}\bar{\phi}\left( \boldsymbol{x}\right) \xi _{n}\left( d%
\boldsymbol{x}\right) ,
\end{equation*}%
where $\bar{\phi}\left( \boldsymbol{x}_{i}\right) =\left( 1/n_{i}\right)
\sum_{j=1}^{n_{i}}\phi \left( Y_{j}\left( \boldsymbol{x}_{i}\right) \right) $%
. We abbreviate $\sum_{\left\{ i|n_{i}>0\right\} }\sum_{j=1}^{n_{i}}$ by $%
\sum_{i,j}$.
\end{remark}

Applying the theorem of \cite{vw98} mentioned in Remark \ref{remark: assumptions} 
requires that we verify further properties of the limit, given in Lemma \ref%
{lemma: limit} below.


\begin{lemma}
\label{lemma: limit} As well as C1) - C3) assume that the parameter space $%
\Theta $ is a compact subset of $\mathbb{R}^{p}$. Then with $\Psi _{n}\left( 
\boldsymbol{\theta }\right) \overset{def}{=}n^{-1}\sum_{i,j}\psi _{\sigma
}\left( Y_{j}\left( \boldsymbol{x}_{i}\right) -\boldsymbol{f}^{\T
}\left( \boldsymbol{x}_{i}\right) \left( \boldsymbol{\theta }+\boldsymbol{%
\tilde{\theta}}_{n}\right) \right) \boldsymbol{f}\left( \boldsymbol{x}%
_{i}\right) ,$ the function 
\begin{equation*}
\Psi \left( \boldsymbol{\theta }\right) =\int_{\mathcal{\chi }}E\left[ \psi
_{\sigma }\left( Y\left( \boldsymbol{x}\right) -\boldsymbol{f}^{\T
}\left( \boldsymbol{x}\right) \left( \boldsymbol{\theta }+\boldsymbol{\theta 
}_{0}\right) \right) \right] \boldsymbol{f}\left( \boldsymbol{x}\right) \xi
_{\ast }\left( d\boldsymbol{x}\right)
\end{equation*}%
satisfies 
\begin{eqnarray*}
(i) &&\sup_{\boldsymbol{\theta }\in \boldsymbol{\Theta }}\left\Vert \Psi _{n}\left( \boldsymbol{%
\theta }\right) -\Psi \left( \boldsymbol{\theta }\right) \right\Vert \overset%
{pr}{\rightarrow }0, \\
(ii) &\text{ }&\text{for every }\delta >0\text{, }\inf_{\left\Vert 
\boldsymbol{\theta }\right\Vert \geq \delta }\left\Vert \Psi \left( 
\boldsymbol{\theta }\right) \right\Vert >0=\left\Vert\Psi ( \boldsymbol{0}_p\right)\Vert
.
\end{eqnarray*}
\end{lemma}


\subsection{Proof of Lemma \protect\ref{lemma: limit}}

With $\boldsymbol{\tilde{\Psi}}_{n}(\boldsymbol{\theta })\overset{def}{=}%
\int_{\chi }E\!\left[ \psi _{\sigma }\!\left( Y(\boldsymbol{x})-\boldsymbol{f%
}^{\T }(\boldsymbol{x})(\boldsymbol{\theta }+\widetilde{\boldsymbol{%
\theta }}_{n})\right) \right] \boldsymbol{f}(\boldsymbol{x})\,\xi _{n}(d%
\boldsymbol{x})=E\left[ \Psi _{n}\left( \boldsymbol{\theta }\right) \right] $
we shall establish (i) of the Lemma by showing:

\begin{enumerate}
\item[(a)] $\sup_{\boldsymbol{\theta }\in \boldsymbol{\Theta }}\left\Vert 
\boldsymbol{\tilde{\Psi}}_{n}\left( \boldsymbol{\theta }\right) -\boldsymbol{%
\Psi }\left( \boldsymbol{\theta }\right) \right\Vert \rightarrow 0$,

\item[(b)] $\sup_{\boldsymbol{\theta }\in \boldsymbol{\Theta }}\left\Vert 
\boldsymbol{\Psi }_{n}\left( \boldsymbol{\theta }\right) -\boldsymbol{\tilde{%
\Psi}}_{n}\left( \boldsymbol{\theta }\right) \right\Vert \overset{pr}{%
\rightarrow }0$.
\end{enumerate}

For (a), with $\boldsymbol{h}(\boldsymbol{x},\boldsymbol{\theta },%
\boldsymbol{\eta })\overset{def}{=}E\!\left[ \psi _{\sigma }\!\left( Y(%
\boldsymbol{x})-\boldsymbol{f}^{\T }(\boldsymbol{x})(\boldsymbol{\theta }%
+\boldsymbol{\eta })\right) \right] \boldsymbol{f}(\boldsymbol{x})$ we have $%
\boldsymbol{\Psi }(\boldsymbol{\theta })=\int_{\chi }\boldsymbol{h}(%
\boldsymbol{x},\boldsymbol{\theta },\boldsymbol{\theta }_{0})\,\xi _{\ast }(d%
\boldsymbol{x})$ and $\boldsymbol{\tilde{\Psi}}_{n}\left( \boldsymbol{\theta 
}\right) =\int_{\chi }\boldsymbol{h}(\boldsymbol{x},\boldsymbol{\theta },%
\boldsymbol{\tilde{\theta}}_{n})\,\xi _{n}(d\boldsymbol{x})$, with 
\begin{eqnarray*}
\boldsymbol{\Psi }(\boldsymbol{\theta })-\boldsymbol{\tilde{\Psi}}_{n}\left( 
\boldsymbol{\theta }\right) &=&\int_{\chi }\left( \boldsymbol{h}(\boldsymbol{%
x},\boldsymbol{\theta },\boldsymbol{\theta }_{0})-\boldsymbol{h}(\boldsymbol{%
x},\boldsymbol{\theta },\boldsymbol{\tilde{\theta}}_{n})\right) \,\xi _{n}(d%
\boldsymbol{x})+\int_{\chi }\boldsymbol{h}(\boldsymbol{x},\boldsymbol{\theta 
},\boldsymbol{\theta }_{0})\,(\xi _{\ast }-\xi _{n})(d\boldsymbol{x}) \\
&=&\boldsymbol{\alpha }_{n}(\boldsymbol{\theta })+\boldsymbol{\beta }_{n}(%
\boldsymbol{\theta })\text{, say.}
\end{eqnarray*}%
Since $\boldsymbol{h}(\boldsymbol{x},\boldsymbol{\theta },\boldsymbol{\eta }%
) $ is differentiable w.r.t. $\boldsymbol{\eta }$, and both $\psi _{\sigma
}^{\prime }$ and $\boldsymbol{f}$ are bounded, there is $C>0$ for which $%
\left\Vert \boldsymbol{h}(\boldsymbol{x},\boldsymbol{\theta },\boldsymbol{%
\theta }_{0})-\boldsymbol{h}(\boldsymbol{x},\boldsymbol{\theta },\boldsymbol{%
\tilde{\theta}}_{n})\right\Vert \leq C\left\Vert \boldsymbol{\theta }_{0}-%
\boldsymbol{\tilde{\theta}}_{n}\right\Vert $, uniformly in $\boldsymbol{x}%
\in \chi $ and $\boldsymbol{\theta }\in \Theta $. Hence 
\begin{equation}
\sup_{\boldsymbol{\theta }\in \Theta }\left\Vert \boldsymbol{\alpha }_{n}(%
\boldsymbol{\theta })\right\Vert \leq C\left\Vert \boldsymbol{\theta }_{0}-%
\boldsymbol{\tilde{\theta}}_{n}\right\Vert \rightarrow 0.  \label{first sup}
\end{equation}%
With $\boldsymbol{g}_{\boldsymbol{\theta }}(\boldsymbol{x})\overset{def}{=}%
\boldsymbol{h}(\boldsymbol{x},\boldsymbol{\theta },\boldsymbol{\theta }_{0})$
we obtain, in a similar fashion, $\sup_{\boldsymbol{x}\in \chi }\Vert 
\boldsymbol{g}_{\boldsymbol{\theta }_{1}}(\boldsymbol{x})-\boldsymbol{g}_{%
\boldsymbol{\theta }_{2}}(\boldsymbol{x})\Vert \leq C\Vert \boldsymbol{%
\theta }_{1}-\boldsymbol{\theta }_{2}\Vert $. Then since $\boldsymbol{\Theta 
}$ is compact, for every $\delta >0$ there exist $\boldsymbol{\theta }%
^{(1)},\dots ,\boldsymbol{\theta }^{(K)}\in \Theta $ such that every $%
\boldsymbol{\theta }\in \boldsymbol{\Theta }$ lies within distance $\delta $
of one of these points. Therefore, 
\begin{equation*}
\sup_{\boldsymbol{\theta }\in \Theta }\left\Vert \boldsymbol{\beta }_{n}(%
\boldsymbol{\theta })\right\Vert \leq 2C\delta +\max_{1\leq k\leq
K}\left\Vert \int_{\chi }\boldsymbol{g}_{\boldsymbol{\theta }^{(k)}}(%
\boldsymbol{x})\,\xi _{\ast }(d\boldsymbol{x})-\int_{\chi }\boldsymbol{g}_{%
\boldsymbol{\theta }^{(k)}}(\boldsymbol{x})\,\xi _{n}(d\boldsymbol{x}%
)\right\Vert .
\end{equation*}%
That $\xi _{n}\overset{d}{\rightarrow }\xi _{\ast }$ implies that the
preceding norms, hence their maximum, converge to $0$. Letting $\delta
\rightarrow 0$, we conclude that $\sup_{\boldsymbol{\theta }\in \Theta
}\left\Vert \boldsymbol{\beta }_{n}(\boldsymbol{\theta })\right\Vert
\rightarrow 0$, which together with (\ref{first sup}) yields (a).

For (b), we write $\boldsymbol{\Delta }_{n}(\boldsymbol{\theta })\overset{def%
}{=}\boldsymbol{\Psi }_{n}(\boldsymbol{\theta })-\boldsymbol{\tilde{\Psi}}%
_{n}\left( \boldsymbol{\theta }\right) \ =n^{-1}\sum_{i,j}\left\{ 
\boldsymbol{\zeta }_{ij,n}(\boldsymbol{\theta })-E[\boldsymbol{\zeta }%
_{ij,n}(\boldsymbol{\theta })]\right\} ,$ where 
\begin{equation*}
\boldsymbol{\zeta }_{ij,n}(\boldsymbol{\theta })=\psi _{\sigma }\!\left(
Y_{j}(\boldsymbol{x}_{i})-\boldsymbol{f}^{\T }\left( \boldsymbol{x}%
_{i}\right) (\boldsymbol{\theta }+\widetilde{\boldsymbol{\theta }}%
_{n})\right) \boldsymbol{f}(\boldsymbol{x}_{i}).
\end{equation*}%
Arguing as above, there exists $C>0$ such that for all $\boldsymbol{\theta }%
_{1},\boldsymbol{\theta }_{2}\in \boldsymbol{\Theta }$ and each pair $(i,j)$%
, 
\begin{equation*}
\left\Vert \boldsymbol{\zeta }_{ij,n}(\boldsymbol{\theta }_{1})-\boldsymbol{%
\zeta }_{ij,n}(\boldsymbol{\theta }_{2})\right\Vert \leq C\left\Vert 
\boldsymbol{\theta }_{1}-\boldsymbol{\theta }_{2}\right\Vert .
\end{equation*}%
Taking expectations gives the same bound for the centred version, so that 
\begin{equation*}
\left\Vert \boldsymbol{\Delta }_{n}(\boldsymbol{\theta }_{1})-\boldsymbol{%
\Delta }_{n}(\boldsymbol{\theta }_{2})\right\Vert \leq 2C\left\Vert 
\boldsymbol{\theta }_{1}-\boldsymbol{\theta }_{2}\right\Vert .
\end{equation*}%
Thus $\Delta _{n}$ is uniformly Lipschitz on $\boldsymbol{\Theta }$, with a
constant independent of $n$.

Fix $\varepsilon >0$. Choose $\delta >0$ such that $2C\delta <\varepsilon /2$%
, and let $\boldsymbol{\theta }^{(1)},\dots ,\boldsymbol{\theta }^{(K)}$ be
a finite $\delta $-net of $\Theta $, as above. Then 
\begin{equation*}
\sup_{\boldsymbol{\theta }\in \Theta }\left\Vert \boldsymbol{\Delta }_{n}(%
\boldsymbol{\theta })\right\Vert \leq \max_{1\leq k\leq K}\left\Vert 
\boldsymbol{\Delta }_{n}(\boldsymbol{\theta }^{(k)})\right\Vert +\varepsilon
/2,
\end{equation*}%
and it remains only to show that 
\begin{equation}
\max_{1\leq k\leq K}\left\Vert \boldsymbol{\Delta }_{n}(\boldsymbol{\theta }%
^{(k)})\right\Vert \overset{pr}{\rightarrow }0.  \label{remains}
\end{equation}%
Write $\Delta _{n}(\boldsymbol{\theta }^{(k)})=n^{-1}\sum_{i,j}\left( 
\boldsymbol{\zeta }_{ij,n}(\boldsymbol{\theta }^{(k)})-E\left[ \boldsymbol{%
\zeta }_{ij,n}(\boldsymbol{\theta }^{(k)})\right] \right) =n^{-1}\sum_{i,j}%
\boldsymbol{Z}_{ij,n}^{(k)}$, say. We can write $\boldsymbol{\zeta }_{ij,n}(%
\boldsymbol{\theta })$ as $\boldsymbol{\zeta }_{ij,n}(\boldsymbol{\theta })=\psi _{\sigma }\!\left(
\varepsilon _{ij}+e_{i,n}\left( \boldsymbol{\theta }\right) \right) 
\boldsymbol{f}(\boldsymbol{x}_{i})
$,
for certain functions $e_{i,n}\left( \boldsymbol{\theta }\right) $,
continuous on $\boldsymbol{\Theta }$ hence uniformly bounded there. Then $E[%
\boldsymbol{Z}_{ij,n}^{(k)}]=\boldsymbol{0}$ and{\ there is }$C>0$ for which 
{$\sup_{i,j,n,k}E\Vert \boldsymbol{Z}_{ij,n}^{(k)}\Vert ^{2}\leq C$}. Since
the $\boldsymbol{Z}_{ij,n}^{(k)}$ are independent for each $(k,n)$, \textsc{%
var}$\left[ \Delta _{n}(\boldsymbol{\theta }^{(k)})\right] \leq
n^{-1}C\rightarrow 0$, implying by Chebyshev's inequality that $\boldsymbol{%
\Delta }_{n}(\boldsymbol{\theta }^{(k)})\overset{pr}{\rightarrow }0$ for
each $k=1,\dots ,K$. Now (\ref{remains}) follows, thus completing the proof
of (b).

For part (ii) of the Lemma, recall that the design space
is finite and that each $\xi _{n}\left( \boldsymbol{x}_{i}\right) $ tends to 
$\xi _{\ast }\left( \boldsymbol{x}_{i}\right) $, so that
\begin{equation}
\int_{\mathcal{\chi }}E\left[ \psi _{\sigma }\left( Y\left( \boldsymbol{x}%
\right) -\boldsymbol{f}^{\T }\left( \boldsymbol{x}\right) \boldsymbol{%
\theta }_{0}\right) \right] \boldsymbol{f}\left( \boldsymbol{x}\right) \xi
_{\ast }\left( d\boldsymbol{x}\right) =\lim_{n\rightarrow \infty }\int_{%
\mathcal{\chi }}E\left[ \psi _{\sigma }\left( \varepsilon +\tau \left( 
\boldsymbol{x}\right) \right) \right] \boldsymbol{f}\left( \boldsymbol{x}%
\right) \xi _{n}\left( d\boldsymbol{x}\right) =\boldsymbol{0}_p,  \label{root}
\end{equation}%
since $\tau \left( \boldsymbol{x}\right) $ is $O\left( n^{-1/2}\right) $. It
follows that $\Psi \left( \boldsymbol{0}_p\right) =\boldsymbol{0}_p $.

By (\ref{root}), $\boldsymbol{\theta }_{0}$ is a stationary point of the
function $k\left( \boldsymbol{\theta }\right) \overset{def}{=}\int_{\mathcal{%
\chi }}E\left[ \rho _{\sigma }\left( Y\left( \boldsymbol{x}\right) -%
\boldsymbol{f}^{\T }\left( \boldsymbol{x}\right) \boldsymbol{\theta }%
\right) \right] \xi _{\ast }\left( d\boldsymbol{x}\right) $, where $\rho
_{\sigma }(x)=\int^{x}\psi _{\sigma }\left( t\right) dt$. The Hessian is $%
\ddot{k}\left( \boldsymbol{\theta }\right) =\int_{\mathcal{\chi }}E\left[
\psi _{\sigma }^{\prime }\left( \varepsilon +M_{\boldsymbol{\theta }}\left( 
\boldsymbol{x}\right) \right) \right] \boldsymbol{f}\left( \boldsymbol{x}%
\right) \boldsymbol{f}^{\T }\left( \boldsymbol{x}\right) \xi _{\ast
}\left( d\boldsymbol{x}\right) ,$ for $M_{\boldsymbol{\theta }}\left( 
\boldsymbol{x}\right) =E\left[  Y\left( \boldsymbol{x}%
\right)] -\boldsymbol{f} ^{\T }\left( \boldsymbol{x}\right) \boldsymbol{\theta }\right) $. Since $\psi_{\sigma }$ is weakly increasing, $\psi _{\sigma }^{\prime }\left(
\varepsilon +M_{\boldsymbol{\theta }}\left( \boldsymbol{x}\right) \right)
\geq 0$ and so $E\left[ \psi _{\sigma }^{\prime }\left( \varepsilon +M_{%
\boldsymbol{\theta }}\left( \boldsymbol{x}\right) \right) \right] $ must be
strictly positive, else $\psi _{\sigma }^{\prime }\left( \varepsilon +M_{%
\boldsymbol{\theta }}\left( \boldsymbol{x}\right) \right) \equiv 0$ (a.s.).
Put 
\begin{equation*}
c\left( \boldsymbol{\theta }\right) =\min_{i=1,...,N}\left\{ E\left[ \psi
_{\sigma }^{\prime }\left( \varepsilon +M_{\boldsymbol{\theta }}\left( 
\boldsymbol{x}_{i}\right) \right) \right] \left\vert {}\right. \xi _{\ast
}\left( \boldsymbol{x}_{i}\right) >0\right\} .
\end{equation*}%
Then $c\left( \boldsymbol{\theta }\right) >0$ and $\ddot{k}\left( 
\boldsymbol{\theta }\right) \succeq c\left( \boldsymbol{\theta }\right)
\int_{\mathcal{\chi }}\boldsymbol{f}\left( \boldsymbol{x}\right) \boldsymbol{%
f}^{\T }\left( \boldsymbol{x}\right) \xi _{\ast }\left( d\boldsymbol{x}%
\right) \succ \boldsymbol{0},$ so that $\boldsymbol{\theta }_{0}$ is the
unique stationary point of the strictly convex function $k\left( \boldsymbol{%
\theta }\right) $ and $\boldsymbol{0}_p$ is the unique zero of $\left\Vert 
\boldsymbol{\Psi }\left( \boldsymbol{\theta }\right) \right\Vert $.

Since $\boldsymbol{\Theta }$ is compact, its closed subset $\boldsymbol{%
\Theta }_{\delta }$ defined by $\left\Vert \boldsymbol{\theta }\right\Vert
\geq \delta $ is also compact, hence the continuous function $\left\Vert 
\boldsymbol{\Psi }\left( \boldsymbol{\theta }\right) \right\Vert $ attains
its $\inf $ there:%
\begin{equation*}
\inf_{\boldsymbol{\Theta }_{\delta }}\left\Vert \boldsymbol{\Psi }\left( 
\boldsymbol{\theta }\right) \right\Vert =\min_{\boldsymbol{\Theta }_{\delta
}}\left\Vert \boldsymbol{\Psi }\left( \boldsymbol{\theta }\right)
\right\Vert >0.
\end{equation*}%
This proves (ii), and completes the proof of Lemma \ref{lemma: limit}%
.\hfill\ $\square $

\medskip

\subsection{Proof of Theorem \protect\ref{thm:main}}
We must show that $\sqrt{n}\boldsymbol{M}_{n}^{1/2}\left( \boldsymbol{\hat{%
\theta}}_{n}-\boldsymbol{\tilde{\theta}}_{n}\right) \overset{d}{%
\longrightarrow }N\left( \boldsymbol{0}_{p},\sigma _{M}^{2}\boldsymbol{I}%
_{p}\right) $. The stochastic component of the estimate will be expressed in terms of 
\begin{equation}
\boldsymbol{\bar{z}}_{n}\overset{def}{=}n^{-1}\sum_{i,j}\psi _{\sigma
}\left( \varepsilon _{ij}\right) \boldsymbol{f}\left( \boldsymbol{x}%
_{i}\right) . \label{znbar}
\end{equation}
Note that $\boldsymbol{\bar{z}}_{n}$ is an average of
independent r.v.s, each with mean zero, and covariance matrices%
\begin{equation*}
\text{\textsc{var}}\left[ \psi _{\sigma }\left( \varepsilon \right) \right] 
\boldsymbol{f}\left( \boldsymbol{x}_{i}\right) \boldsymbol{f}^{{\mathrm{%
\scriptscriptstyle T}}}\left( \boldsymbol{x}_{i}\right) =E\left[ \psi
_{\sigma }^{2}\left( \varepsilon \right) \right] \boldsymbol{f}\left( 
\boldsymbol{x}_{i}\right) \boldsymbol{f}^{{\mathrm{\scriptscriptstyle T}}%
}\left( \boldsymbol{x}_{i}\right) ,
\end{equation*}%
so that \textsc{cov}$\left[ \boldsymbol{\bar{z}}_{n}\right] =n^{-1}E\left[
\psi _{\sigma }^{2}\left( \varepsilon \right) \right] \boldsymbol{M}%
_{n}\succ \boldsymbol{0}_{p}$. By the Lindberg-Feller CLT, 
\begin{equation}
\boldsymbol{M}_{n}^{-1/2}\sqrt{n}\boldsymbol{\bar{z}}_{n}\overset{d}{%
\rightarrow }N\left( \boldsymbol{0}_{p},E\left[ \psi _{\sigma }^{2}\left(
\varepsilon \right) \right] \boldsymbol{I}_{p}\right) .  \label{AN-z}
\end{equation}

Completing the proof of Theorem \ref{thm:main} requires a preliminary Lemma. 

\begin{lemma}
\label{root-n consistent} Recall that $\boldsymbol{\tilde{\theta}}_{n}=%
\boldsymbol{\theta }_{0}+\boldsymbol{M}_{n}^{-1}\boldsymbol{b}_{n}$ . We
have that $\boldsymbol{\hat{\theta}}_{n}-\boldsymbol{\tilde{\theta}}_{n}$ is
a $\sqrt{n}$-consistent estimate of $\boldsymbol{0}_{p}$:    
\begin{equation}
\boldsymbol{\hat{\theta}}_{n}-\boldsymbol{\tilde{\theta}}_{n}=O_{p}\left(
n^{-1/2}\right) .  \label{claim2}
\end{equation}
\end{lemma}

\noindent \textit{Proof}: We use the functions 
\begin{eqnarray*}
a_{n}(\boldsymbol{x}) &=&\boldsymbol{f}^{\T}\left( \boldsymbol{x}\right) \left( \boldsymbol{\tilde{\theta}}_{n}-%
\boldsymbol{\theta }_{0}\right) =\boldsymbol{f}^{\T}\left( \boldsymbol{x}%
\right) \boldsymbol{M}_{n}^{-1}\boldsymbol{b}_{n}, \\
h_{n}\left( \boldsymbol{x}\right)  &=&\tau \left( \boldsymbol{x}\right)
-a_{n}(\boldsymbol{x}), \\
d_{n}(\boldsymbol{x}) &=&\boldsymbol{f}^{{\mathrm{\scriptscriptstyle T}}%
}\left( \boldsymbol{x}\right) \left( \boldsymbol{\hat{\theta}}_{n}-%
\boldsymbol{\tilde{\theta}}_{n}\right) .
\end{eqnarray*}%
By (\ref{bound}) and C2), both $a_{n}(\boldsymbol{x})$ and $h_{n}\left( \boldsymbol{x%
}\right) $ are $o_{p}\left( 1\right) $.

We first show that $\boldsymbol{\hat{\theta}}_{n}-\boldsymbol{\tilde{\theta}}%
_{n}$ is a consistent estimate of $\boldsymbol{0}_{p}$: 
\begin{equation}
\boldsymbol{\hat{\theta}}_{n}-\boldsymbol{\tilde{\theta}}_{n}=o_{p}\left(
1\right) .  \label{claim1}
\end{equation}%
We have that $\Psi _{n}\left( \boldsymbol{\hat{\theta}}_{n}-\boldsymbol{%
\tilde{\theta}}_{n}\right) =n^{-1}\sum_{i,j}\psi _{\sigma }\left(
Y_{j}\left( \boldsymbol{x}_{i}\right) -\boldsymbol{f}^{\T}\left( \boldsymbol{x%
}_{i}\right) \boldsymbol{\hat{\theta}}_{n}\right) \boldsymbol{f}\left( 
\boldsymbol{x}_{i}\right) =\boldsymbol{0}_{p}$, by (\ref{m-eqn}). Now
Theorem 5.9 of \cite{vw98}, with $\boldsymbol{\theta }_{0}=$ $\boldsymbol{0}%
_{p}$ and under the conditions there -- ensured by our Lemma \ref{lemma:
limit} -- asserts (\ref{claim1}), and also that $d_{n}(\boldsymbol{x}%
)=o_{p}\left( 1\right) $.

We require the stronger $\sqrt{n}$-consistency. For this, observe that by (\ref{m-eqn}), 
\begin{eqnarray}
\boldsymbol{0}_{p} &=&\frac{1}{n}\sum_{i,j}\psi _{\sigma }\left( Y_{j}\left( 
\boldsymbol{x}_{i}\right) -\boldsymbol{f}^{{\mathrm{\scriptscriptstyle T}}%
}\left( \boldsymbol{x}_{i}\right) \boldsymbol{\hat{\theta}}_{n}\right) 
\boldsymbol{f}\left( \boldsymbol{x}_{i}\right)   \notag \\
&=&\frac{1}{n}\sum_{i,j}\psi _{\sigma }\left( Y_{j}\left( \boldsymbol{x}%
_{i}\right) -\boldsymbol{f}^{{\mathrm{\scriptscriptstyle T}}}\left( 
\boldsymbol{x}_{i}\right) \boldsymbol{\tilde{\theta}}_{n}-d_{n}(\boldsymbol{x%
}_{i})\right) \boldsymbol{f}\left( \boldsymbol{x}_{i}\right)   \notag \\
&=&\frac{1}{n}\sum_{i,j}\psi _{\sigma }\left( Y_{j}\left( \boldsymbol{x}%
_{i}\right) -\boldsymbol{f}^{{\mathrm{\scriptscriptstyle T}}}\left( 
\boldsymbol{x}_{i}\right) \boldsymbol{\tilde{\theta}}_{n}\right) \boldsymbol{%
f}\left( \boldsymbol{x}_{i}\right) -\boldsymbol{G}_{n}\left( \boldsymbol{%
\hat{\theta}}_{n}-\boldsymbol{\tilde{\theta}}_{n}\right) ,  \label{exp1}
\end{eqnarray}%
where for generic $t_{i}\in \lbrack 0,1]$ and $u_{ij}\in \lbrack 0,1]$, and
with $s_{n,i}=h_{n}(\boldsymbol{x}_{i})-t_{i}d_{n}(\boldsymbol{x}%
_{i})=o_{p}\left( 1\right) $, 
\begin{eqnarray*}
\boldsymbol{G}_{n} &=&\frac{1}{n}\sum_{i,j}\psi _{\sigma }^{\prime }\left(
Y_{j}\left( \boldsymbol{x}_{i}\right) -\boldsymbol{f}^{\T}\left( \boldsymbol{x%
}_{i}\right) \boldsymbol{\tilde{\theta}}_{n}-t_{i}d_{n}(\boldsymbol{x}%
_{i})\right) \boldsymbol{f}\left( \boldsymbol{x}_{i}\right) \boldsymbol{f}^{{%
\mathrm{\scriptscriptstyle T}}}\left( \boldsymbol{x}_{i}\right)  \\
&=&\frac{1}{n}\sum_{i,j}\psi _{\sigma }^{\prime }\left( \varepsilon
_{ij}+s_{n,i}\right) \boldsymbol{f}\left( \boldsymbol{x}_{i}\right) 
\boldsymbol{f}^{\T}\left( \boldsymbol{x}%
_{i}\right)  \\
&=&\frac{1}{n}\sum_{i,j}\psi _{\sigma }^{\prime }\left( \varepsilon
_{ij}\right) \boldsymbol{f}\left( \boldsymbol{x}_{i}\right) \boldsymbol{f}^{{%
\mathrm{\scriptscriptstyle T}}}\left( \boldsymbol{x}_{i}\right) +\frac{1}{n}%
\sum_{i,j}\psi _{\sigma }^{\prime \prime }\left( \varepsilon
_{ij}+u_{ij}s_{n,i}\right) s_{n,i}\boldsymbol{f}\left( \boldsymbol{x}%
_{i}\right) \boldsymbol{f}^{{\mathrm{\scriptscriptstyle T}}}\left( 
\boldsymbol{x}_{i}\right) .
\end{eqnarray*}%
We will show that $\boldsymbol{G}_{n}$ tends in probability to the positive
definite matrix $\boldsymbol{G}_{\ast }=E\left[ \psi _{\sigma }^{\prime
}\left( \varepsilon \right) \right] \int_{\mathcal{\chi }}\boldsymbol{f}%
\left( \boldsymbol{x}\right) \boldsymbol{f}^{\T}\left( \boldsymbol{x}\right)
\xi _{\ast }\left( d\boldsymbol{x}\right) $. Then by (\ref{exp1}), and ignoring terms which are $o_p(n^{-1/2}),$%
\begin{eqnarray*}
\sqrt{n}\left( \boldsymbol{\hat{\theta}}_{n}-\boldsymbol{\tilde{\theta}}%
_{n}\right)  &=&\left( \boldsymbol{G}_{\ast }^{-1}+o_{p}(1)\right) \left\{ 
\frac{1}{\sqrt{n}}\sum_{i,j}\psi _{\sigma }\left( Y_{j}\left( \boldsymbol{x}%
_{i}\right) -\boldsymbol{f}^{{\mathrm{\scriptscriptstyle T}}}\left( 
\boldsymbol{x}_{i}\right) \boldsymbol{\tilde{\theta}}_{n}\right) \boldsymbol{%
f}\left( \boldsymbol{x}_{i}\right) \right\}  \\
&=&\left( \boldsymbol{G}_{\ast }^{-1}+o_{p}(1)\right) \left\{ \frac{1}{\sqrt{%
n}}\sum_{i,j}\psi _{\sigma }\left( \varepsilon _{ij}+h_{n}\left( \boldsymbol{%
x}_{i}\right) \right) \boldsymbol{f}\left( \boldsymbol{x}_{i}\right)
\right\}  \\
&=&\left( \boldsymbol{G}_{\ast }^{-1}+o_{p}(1)\right) \left\{ \sqrt{n}%
\boldsymbol{\bar{z}}_{n}+\frac{1}{n}\sum_{i,j}\psi _{\sigma }^{\prime
}\left( \varepsilon _{ij}+t_{i}h_{n}\left( \boldsymbol{x}_{i}\right) \right) 
\sqrt{n}h\left( \boldsymbol{x}_{i}\right) \boldsymbol{f}\left( \boldsymbol{x}%
_{i}\right) \right\} ,
\end{eqnarray*}%
where the first term in the final set of braces is $O_{p}(1)$ by (\ref{AN-z}), and the second
is $O_{p}(1)$ because 
\begin{equation*}
\sqrt{n}h_{n}\left( \boldsymbol{x}\right) =\sqrt{n}\tau \left( \boldsymbol{x}%
\right) -\sqrt{n}\boldsymbol{f}^{{\mathrm{\scriptscriptstyle T}}}\left( 
\boldsymbol{x}\right) \boldsymbol{M}_{n}^{-1}\boldsymbol{b}_{n}=\sqrt{n}%
\tau \left( \boldsymbol{x}\right) -\boldsymbol{f}^{{\mathrm{%
\scriptscriptstyle T}}}\left( \boldsymbol{x}\right) \boldsymbol{M}%
_{n}^{-1}\int_{\mathcal{\chi }}\boldsymbol{f}\left( \boldsymbol{x}\right) 
\sqrt{n}\tau \left( \boldsymbol{x}\right) \xi _{n}\left( d\boldsymbol{x}%
\right) 
\end{equation*}%
is bounded, again by virtue of (\ref{bound}).\bigskip  Thus showing that \begin{equation}
\boldsymbol{G}_{n}\overset{pr}{\rightarrow }\boldsymbol{G}_{\ast },
\label{A-limit}
\end{equation}%
will complete the verification of (\ref{claim2}).

To show (\ref{A-limit}) note that by the triangle inequality,
\begin{equation}
\left\Vert \boldsymbol{G}_{n}-\boldsymbol{G}_{\ast }\right\Vert \leq
\left\Vert \boldsymbol{G}_{n}-\frac{1}{n}\sum_{i,j}\psi _{\sigma }^{\prime
}\left( \varepsilon _{ij}\right) \boldsymbol{f}\left( \boldsymbol{x}%
_{i}\right) \boldsymbol{f}^{{\mathrm{\scriptscriptstyle T}}}\left( 
\boldsymbol{x}_{i}\right) \right\Vert +\left\Vert \boldsymbol{W}%
_{n}\right\Vert +\left\Vert \frac{1}{n}\sum_{i,j}E\left[ \psi _{\sigma
}^{\prime }\left( \varepsilon _{ij}\right) \right] \boldsymbol{f}\left( 
\boldsymbol{x}_{i}\right) \boldsymbol{f}^{{\mathrm{\scriptscriptstyle T}}%
}\left( \boldsymbol{x}_{i}\right) -\boldsymbol{G}_{\ast }\right\Vert ,
\label{A-decomp}
\end{equation}%
where 
\begin{equation*}
\boldsymbol{W}_{n}\overset{def}{=}\frac{1}{n}\sum_{i,j}\left( \psi _{\sigma
}^{\prime }\left( \varepsilon _{ij}\right) -E\left[ \psi _{\sigma }^{\prime
}\left( \varepsilon _{ij}\right) \right] \right) \boldsymbol{f}\left( 
\boldsymbol{x}_{i}\right) \boldsymbol{f}^{{\mathrm{\scriptscriptstyle T}}%
}\left( \boldsymbol{x}_{i}\right) =\sum_{i}\xi _{n}\left( \boldsymbol{x}%
_{i}\right) V_{i}\boldsymbol{f}\left( \boldsymbol{x}_{i}\right) \boldsymbol{f%
}^{{\mathrm{\scriptscriptstyle T}}}\left( \boldsymbol{x}_{i}\right) ,
\end{equation*}%
for independent r.v.s $V_{i}=\sum_{j=1}^{n_{i}}\left( \psi _{\sigma
}^{\prime }\left( \varepsilon _{ij}\right) -E\left[ \psi _{\sigma }^{\prime
}\left( \varepsilon _{ij}\right) \right] \right) /n_{i}$ with means $0$ and
variances \textsc{var}$\left[ \psi _{\sigma }^{\prime }\left( \varepsilon
\right) \right] /n_{i}\leq m_{1}^{2}/n_{i}$.

The first norm on the rhs of (\ref{A-decomp}) is 
\begin{equation*}
\left\Vert \frac{1}{n}\sum_{i,j}\psi _{\sigma }^{\prime \prime }\left(
\varepsilon _{ij}+u_{ij}s_{n,i}\right) s_{n,i}\boldsymbol{f}\left( 
\boldsymbol{x}_{i}\right) \boldsymbol{f}^{{\mathrm{\scriptscriptstyle T}}%
}\left( \boldsymbol{x}_{i}\right) \right\Vert \leq \frac{m_{2}}{n}%
\sum_{i,j}\left\vert s_{n,i}\right\vert \left\Vert \boldsymbol{f}\left( 
\boldsymbol{x}_{i}\right) \right\Vert ^{2}=o_{p}\left( 1\right) ,
\end{equation*}%
and the third is $\left\Vert \int_{\mathcal{\chi }}E\left[ \psi _{\sigma
}^{\prime }\left( \varepsilon \right) \right] \boldsymbol{f}\left( 
\boldsymbol{x}\right) \boldsymbol{f}^{{\mathrm{\scriptscriptstyle T}}}\left( 
\boldsymbol{x}\right) \left( \xi _{n}-\xi _{\ast }\right) \left( d%
\boldsymbol{x}\right) \right\Vert =o(1)$.  A typical element of $W_{n}$ is $%
W_{n,k,l}=\sum_{i}\xi _{n}\left( \boldsymbol{x}_{i}\right) V_{i}f_{k}\left( 
\boldsymbol{x}_{i}\right) f_{l}\left( \boldsymbol{x}_{i}\right) $, with mean
zero and 
\begin{equation*}
\text{\textsc{var}}\left[ W_{n,k,l}\right] =\sum_{i}\frac{n_{i}^{2}}{n^{2}}%
\text{\textsc{var}}\left[ V_{i}\right] \left( f_{k}\left( \boldsymbol{x}%
_{i}\right) f_{l}\left( \boldsymbol{x}_{i}\right) \right) ^{2}\leq \frac{%
m_{1}^{2}}{n}\int_{\mathcal{\chi }}\left( f_{k}\left( \boldsymbol{x}%
_{i}\right) f_{l}\left( \boldsymbol{x}_{i}\right) \right) ^{2}\xi _{n}\left(
d\boldsymbol{x}\right) =O(n^{-1}),
\end{equation*}%
so that by Chebyshev's Inequality, $W_{n,k,l}\overset{pr}{\rightarrow }0$
and so $\boldsymbol{W}_{n}\overset{pr}{\rightarrow }\boldsymbol{0}.$ Now (%
\ref{A-limit}) follows, completing the proof of the Lemma.\hfill\ $\square
\medskip $

We can now complete the proof of Theorem \ref{thm:main}. To establish the asymptotic normality of $\boldsymbol{\hat{\theta}}_{n}-%
\boldsymbol{\tilde{\theta}}_{n}$, we observe that by (\ref{m-eqn}), 
\begin{eqnarray}
\boldsymbol{0}_p &=& \frac{1}{n}\sum_{i,j}\psi _{\sigma }\left( Y_{j}\left( \boldsymbol{x}%
_{i}\right) -\boldsymbol{f}^{\T }\left( \boldsymbol{x}_{i}\right) 
\boldsymbol{\hat{\theta}}_{n}\right) \boldsymbol{f}\left( \boldsymbol{x}%
_{i}\right)  \notag \\
&=&\frac{1}{n}\sum_{i,j}\psi _{\sigma }\left( Y_{j}\left( 
\boldsymbol{x}_{i}\right) -\boldsymbol{f}^{{\mathrm{\scriptscriptstyle T}}%
}\left( \boldsymbol{x}_{i}\right) \boldsymbol{\tilde{\theta}}%
_{n}-d_{n}\left( \boldsymbol{x}_{i}\right) \right) \boldsymbol{f}\left( 
\boldsymbol{x}_{i}\right)  \notag \\
\text{ }& =&\frac{1}{n}\sum_{i,j}\psi _{\sigma }\left( Y_{j}\left( \boldsymbol{x}%
_{i}\right) -\boldsymbol{f}^{{\mathrm{\scriptscriptstyle T}}}\left( 
\boldsymbol{x}_{i}\right) \boldsymbol{\tilde{\theta}}_{n}\right) \boldsymbol{%
f}\left( \boldsymbol{x}_{i}\right)   \notag \\
\text{}&&-\frac{1}{n}\sum_{i,j}\psi _{\sigma }^{\prime }\left( Y_{j}\left( \boldsymbol{x}%
_{i}\right) -\boldsymbol{f}^{{\mathrm{\scriptscriptstyle T}}}\left( 
\boldsymbol{x}_{i}\right) \boldsymbol{\tilde{\theta}}_{n}\right) \boldsymbol{%
f}\left( \boldsymbol{x}_{i}\right) d_{n}\left( \boldsymbol{x}_{i}\right) +%
\boldsymbol{S}_{n},  \label{exp2}
\end{eqnarray}%
where
\begin{equation*}
\boldsymbol{S}_{n}=\left( 1/2n\right) \sum_{i,j}\psi _{\sigma }^{\prime
\prime }\left( Y_{j}\left( \boldsymbol{x}_{i}\right) -\boldsymbol{f}^{{%
\mathrm{\scriptscriptstyle T}}}\left( \boldsymbol{x}_{i}\right) \boldsymbol{%
\tilde{\theta}}_{n}+t_id_{n}\left( \boldsymbol{x}_{i}\right) \right)
d_{n}^{2}\left( \boldsymbol{x}_{i}\right) \boldsymbol{f}\left( \boldsymbol{x}%
_{i}\right) .
\end{equation*}%
With $B_{n}\overset{def}{=}ch_{\max }\int_{\mathcal{\chi }}\boldsymbol{f}%
\left( \boldsymbol{x}\right) \boldsymbol{f}^{{\mathrm{\scriptscriptstyle T}}%
}\left( \boldsymbol{x}\right) \left\Vert \boldsymbol{f}\left( \boldsymbol{x}%
\right) \right\Vert \xi _{n}\left( d\boldsymbol{x}\right) $ $\leq $ $\int_{%
\mathcal{\chi }}\left\Vert \boldsymbol{f}\left( \boldsymbol{x}\right)
\right\Vert ^{3}\xi _{n}\left( d\boldsymbol{x}\right) <\infty $, we have 
\begin{equation*}
\left\Vert \boldsymbol{S}_{n}\right\Vert \leq \frac{m_{2}}{2n}%
\sum_{i,j}\left( \boldsymbol{\hat{\theta}}_{n}-\boldsymbol{\tilde{\theta}}%
_{n}\right) ^{{\mathrm{\scriptscriptstyle T}}}\boldsymbol{f}\left( 
\boldsymbol{x}_{i}\right) \boldsymbol{f}^{{\mathrm{\scriptscriptstyle T}}%
}\left( \boldsymbol{x}_{i}\right) \left( \boldsymbol{\hat{\theta}}_{n}-%
\boldsymbol{\tilde{\theta}}_{n}\right) \left\Vert \boldsymbol{f}\left( 
\boldsymbol{x}_{i}\right) \right\Vert \leq \frac{m_{2}}{2}B_{n}\left\Vert 
\boldsymbol{\hat{\theta}}_{n}-\boldsymbol{\tilde{\theta}}_{n}\right\Vert
^{2}=O_{p}\left( n^{-1}\right) .
\end{equation*}%
The final step in the preceding line uses Lemma  \ref{root-n consistent}. Thus, from (\ref{exp2}) and disregarding
terms that are $O_{p}\left( n^{-1/2}\right) $, 
\begin{equation}
\sqrt{n}\left( \boldsymbol{\hat{\theta}}_{n}-\boldsymbol{\tilde{\theta}}%
_{n}\right) =\boldsymbol{P}_{n}^{-1}\frac{1}{\sqrt{n}}\sum_{i,j}\psi
_{\sigma }\left( Y_{j}\left( \boldsymbol{x}_{i}\right) -\boldsymbol{f}^{{%
\mathrm{\scriptscriptstyle T}}}\left( \boldsymbol{x}_{i}\right) \boldsymbol{%
\tilde{\theta}}_{n}\right) \boldsymbol{f}\left( \boldsymbol{x}_{i}\right) ,
\label{exp3}
\end{equation}%
where%
\begin{eqnarray}
\boldsymbol{P}_{n} &\overset{def}{=}&\frac{1}{n}\sum_{i,j}\psi _{\sigma
}^{\prime }\left( Y_{j}\left( \boldsymbol{x}_{i}\right) -\boldsymbol{f}^{{%
\mathrm{\scriptscriptstyle T}}}\left( \boldsymbol{x}_{i}\right) \boldsymbol{%
\tilde{\theta}}_{n}\right) \boldsymbol{f}\left( \boldsymbol{x}_{i}\right) 
\boldsymbol{f}^{{\mathrm{\scriptscriptstyle T}}}\left( \boldsymbol{x}%
_{i}\right)   \notag \\
&=&\frac{1}{n}\sum_{i,j}\psi _{\sigma }^{\prime }\left( \tau \left( 
\boldsymbol{x}_{i}\right) +\varepsilon _{ij}-a_{n}\left( \boldsymbol{x}%
_{i}\right) \right) \boldsymbol{f}\left( \boldsymbol{x}_{i}\right) 
\boldsymbol{f}^{{\mathrm{\scriptscriptstyle T}}}\left( \boldsymbol{x}%
_{i}\right)   \notag \\
&=&\frac{1}{n}\sum_{i,j}\psi _{\sigma }^{\prime }\left( \tau \left( 
\boldsymbol{x}_{i}\right) +\varepsilon _{ij}\right) \boldsymbol{f}\left( 
\boldsymbol{x}_{i}\right) \boldsymbol{f}^{{\mathrm{\scriptscriptstyle T}}%
}\left( \boldsymbol{x}_{i}\right) -\boldsymbol{T}_{n},  \label{exp4}
\end{eqnarray}%
and
\begin{equation*}
\boldsymbol{T}_{n}=\frac{1}{n}\sum_{i,j}\psi _{\sigma }^{\prime \prime
}\left( \tau \left( \boldsymbol{x}_{i}\right) +\varepsilon
_{ij}-t_ia_{n}\left( \boldsymbol{x}_{i}\right) \right) a_{n}\left( \boldsymbol{%
x}_{i}\right) \boldsymbol{f}\left( \boldsymbol{x}_{i}\right) \boldsymbol{f}^{%
{\mathrm{\scriptscriptstyle T}}}\left( \boldsymbol{x}_{i}\right) .
\end{equation*}%
For any vector $\boldsymbol{\alpha }$ of unit norm, 
\begin{eqnarray*}
\left\vert \boldsymbol{\alpha }^{{\mathrm{\scriptscriptstyle T}}}\boldsymbol{%
T}_{n}\boldsymbol{\alpha }\right\vert  &\leq &\frac{1}{n}\sum_{i,j}\left%
\vert \psi _{\sigma }^{\prime \prime }\left( \tau \left( \boldsymbol{x}%
_{i}\right) +\varepsilon _{ij}-ta_{n}\left( \boldsymbol{x}_{i}\right)
\right) a_{n}\left( \boldsymbol{x}_{i}\right) \right\vert \left( \boldsymbol{%
\alpha }^{{\mathrm{\scriptscriptstyle T}}}\boldsymbol{f}\left( \boldsymbol{x}%
_{i}\right) \right) ^{2} \\
&\leq &\frac{m_{2}}{n}\sum_{i,j}\left( \boldsymbol{\alpha }^{{\mathrm{%
\scriptscriptstyle T}}}\boldsymbol{f}\left( \boldsymbol{x}_{i}\right)
\right) ^{2}\left( \max_{\mathcal{\chi }}\left\vert a_{n}\left( \boldsymbol{x%
}\right) \right\vert \right)  \\
&=&m_{2}\cdot \boldsymbol{\alpha }^{{\mathrm{\scriptscriptstyle T}}}%
\boldsymbol{M}_{n}\boldsymbol{\alpha }\cdot \max_{\mathcal{\chi }%
}\left\Vert \boldsymbol{f}\left( \boldsymbol{x}\right) \right\Vert \cdot
\left\Vert \boldsymbol{M}_{n}^{-1}\boldsymbol{b}_{n}\right\Vert ,
\end{eqnarray*}%
which is $O\left( n^{-1/2}\right) $. Thus all eigenvalues of $\boldsymbol{T}%
_{n}$ are $O\left( n^{-1/2}\right) $, hence so is $\boldsymbol{T}_{n}$
itself and (\ref{exp4}) becomes 
\begin{equation*}
\boldsymbol{P}_{n}=E\left[ \psi
_{\sigma }^{\prime }\left( \varepsilon \right) \right] \boldsymbol{M}%
_{n}-O_p\left( n^{-1/2}\right) ,
\end{equation*}%
where we invoke the WLLN. Now (\ref{exp3})
becomes%
\begin{equation}
\sqrt{n}\left( \boldsymbol{\hat{\theta}}_{n}-\boldsymbol{\tilde{\theta}}%
_{n}\right) =\left[ E\left[ \psi _{\sigma }^{\prime }\left( \varepsilon
\right) \right] \boldsymbol{M}_{n}+O\left( n^{-1/2}\right) \right]
^{-1}\cdot \sqrt{n}\boldsymbol{\tilde{z}}_{n},  \label{exp5}
\end{equation}%
for%
\begin{equation*}
\boldsymbol{\tilde{z}}_{n}=\frac{1}{n}\sum_{i,j}\psi _{\sigma }\left(
Y_{j}\left( \boldsymbol{x}_{i}\right) -\boldsymbol{f}^{{\mathrm{%
\scriptscriptstyle T}}}\left( \boldsymbol{x}_{i}\right) \boldsymbol{\tilde{%
\theta}}_{n}\right) \boldsymbol{f}\left( \boldsymbol{x}_{i}\right) =\frac{1}{%
n}\sum_{i,j}\psi _{\sigma }\left( \varepsilon _{ij}+h_{n}\left( \boldsymbol{x%
}_{i}\right) \right) \boldsymbol{f}\left( \boldsymbol{x}_{i}\right) .
\end{equation*}%
Recalling that $h_{n}\left( \boldsymbol{x}\right) $ is $O\left(
n^{-1/2}\right) $, and using (\ref{znbar}), we have%
\begin{equation*}
\boldsymbol{\tilde{z}}_{n}=\frac{1}{n}\sum_{i,j}\psi _{\sigma }\left(
\varepsilon _{ij}+h_{n}\left( \boldsymbol{x}_{i}\right) \right) \boldsymbol{f%
}\left( \boldsymbol{x}_{i}\right) =\boldsymbol{\bar{z}}_{n}+\frac{1}{n}%
\sum_{i,j}\psi _{\sigma }^{\prime }\left( \varepsilon _{ij}\right)
h_{n}\left( \boldsymbol{x}_{i}\right) \boldsymbol{f}\left( \boldsymbol{x}%
_{i}\right) +O_{p}\left( n^{-1}\right) ,
\end{equation*}%
with (and again invoking the WLLN)%
\begin{eqnarray*}
\sqrt{n}\boldsymbol{\tilde{z}}_{n}-\sqrt{n}\boldsymbol{\bar{z}}_{n} &=&\frac{%
1}{n}\sum_{i,j}\psi _{\sigma }^{\prime }\left( \varepsilon _{ij}\right) 
\sqrt{n}h_{n}\left( \boldsymbol{x}_{i}\right) \boldsymbol{f}\left( 
\boldsymbol{x}_{i}\right) +O_{p}\left( n^{-1/2}\right)  \\
&=&E\left[ \psi _{\sigma }^{\prime }\left( \varepsilon \right) \right] \int_{%
\mathcal{\chi }}\sqrt{n}h_{n}\left( \boldsymbol{x}\right) \boldsymbol{f}%
\left( \boldsymbol{x}\right) \xi _{n}\left( d\boldsymbol{x}\right)
+o_{p}\left( 1\right)  \\
&=&E\left[ \psi _{\sigma }^{\prime }\left( \varepsilon \right) \right]
\left\{ \sqrt{n}\boldsymbol{b}_{n}-\sqrt{n}\boldsymbol{M}_{n}\left( 
\boldsymbol{M}_{n}^{-1}\boldsymbol{b}_{n}\right) \right\} +o_{p}\left(
1\right) ,
\end{eqnarray*}%
which is $o_{p}\left( 1\right) $. Thus $\sqrt{n}\boldsymbol{\tilde{z}}_{n}=%
\sqrt{n}\boldsymbol{\bar{z}}_{n}+o_{p}(1)$ and (\ref{exp5}) becomes%
\begin{equation*}
\sqrt{n}\left( \boldsymbol{\hat{\theta}}_{n}-\boldsymbol{\tilde{\theta}}%
_{n}\right) =\left[ E\left[ \psi _{\sigma }^{\prime }\left( \varepsilon
\right) \right] \boldsymbol{M}_{n}+O\left( n^{-1/2}\right) \right]
^{-1}\left( \sqrt{n}\boldsymbol{\bar{z}}_{n}+o_{p}\left( 1\right) \right) ,
\end{equation*}%
whence
\begin{equation}
\sqrt{n}\boldsymbol{M}_{n}^{1/2}\left( \boldsymbol{\hat{%
\theta}}_{n}-\boldsymbol{\tilde{\theta}}_{n}\right) =\frac{1}{E\left[ \psi
_{\sigma }^{{\prime} }\left( \varepsilon \right) %
\right] }\boldsymbol{M}_{n}^{-1/2}\sqrt{n}\boldsymbol{\bar{z}}%
_{n}+o_{p}\left( 1\right) .  \label{represent}
\end{equation}%
From this, the statement of the Theorem is immediate.\hfill\ $\square $

\subsection{Proof of Theorem \protect\ref{thm: max loss}}

To maximize $\textsc{aimse}$ over $\tau \in \Upsilon $, note that both of its
first two terms at (\ref{IMSEn}) become larger if $\tau \left( \boldsymbol{x}\right) $
is multiplied by a constant exceeding one in absolute value, hence at a
maximum (\ref{bound}) is attained with equality. Define $
\tau _{0}\left( \boldsymbol{x}\right) =\sqrt{n}\tau \left( \boldsymbol{x}\right) /\kappa \text{ and } \boldsymbol{c} \left( \xi \right) = \sqrt{n} \boldsymbol{b}_0 \left( \xi \right) / \kappa.
$
Then with $\boldsymbol{c}_{n}=\boldsymbol{c}\left( \xi _{n}\right) $ we have that $\max_{\tau \in \Upsilon } \textsc{aimse} \left( \xi _{n}|\tau \right) $  is
\begin{equation}
J\left( \xi _{n}\right) \overset{%
def}{=}\sigma _{M}^{2}tr\boldsymbol{AM}_{n}^{-1}+\kappa
^{2}\left( 1+\max_{\tau _{0}}\boldsymbol{c}_{n}^{\T }\boldsymbol{M}%
_{n}^{-1}\mathbf{A}\boldsymbol{M}_{n}^{-1}\boldsymbol{c}_{n}\right) ,
\label{I1}
\end{equation}%
with $\tau _{0}$ -- which appears in $\boldsymbol{c}_{n}$ via 
$\boldsymbol{b}_{0}$ -- constrained by (\ref{orthogonality}) and by $%
\sum_{i=1}^{N}\tau _{0}^{2}\left( \boldsymbol{x}_{i}\right) =1$.

To carry out the maximization in (\ref{I1}) recall that a feature of the
Gram-Schmidt process used in the construction of $\boldsymbol{Q}$ is that $%
\boldsymbol{F}=\boldsymbol{QT}$ for a non-singular lower triangular matrix $%
\boldsymbol{T}$. Let $\boldsymbol{Q}_{\boldsymbol{\perp }}:N\times N-p$ be
the orthogonal complement of $\boldsymbol{Q}$, so that $\left( \boldsymbol{Q}%
\vdots \boldsymbol{Q}_{\boldsymbol{\perp }}\right) :N\times N$ is
orthogonal. The condition in (\ref{orthogonality}) asserts that $\boldsymbol{\tau }%
_{0}$ is orthogonal to the columns of $\boldsymbol{Q}$, hence is a linear
combination of the columns of $\boldsymbol{Q}_{\boldsymbol{\perp }}$ and so
there is $\boldsymbol{\beta }_{N-p\times 1}$ of unit norm for which $%
\boldsymbol{\tau }_{0}=\boldsymbol{Q}_{\boldsymbol{\perp }}\boldsymbol{\beta 
}$.

In this notation $\boldsymbol{A}=\boldsymbol{F}^{\T }\boldsymbol{F}=%
\boldsymbol{T}^{\T }\boldsymbol{T}$ and 
\begin{eqnarray*}
\boldsymbol{M}_{n} &=&\boldsymbol{F}^{\T }\boldsymbol{D}\left( \xi
_{n}\right) \boldsymbol{F}=\boldsymbol{T}^{\T }\boldsymbol{Q}^{\T }%
\boldsymbol{D}\left( \xi _{n}\right) \boldsymbol{QT}, \\
\boldsymbol{c}_{n} &=&\boldsymbol{F}^{\T }\boldsymbol{D}\left( \xi
_{n}\right) \boldsymbol{\tau }_{0}=\boldsymbol{T}^{\T }\boldsymbol{Q}%
^{\T }\boldsymbol{D}\left( \xi _{n}\right) \boldsymbol{Q}_{\boldsymbol{%
\perp }}{\boldsymbol{\beta }};
\end{eqnarray*}%
then (\ref{I1}) becomes%
\begin{equation}
J\left( \xi _{n}\right) =\sigma _{M}^{2}tr\left( \boldsymbol{Q}^{\T }%
\boldsymbol{D}\left( \xi _{n}\right) \boldsymbol{Q}\right) ^{-1}+\kappa
^{2}\left( 1+\max_{\left\Vert \boldsymbol{\beta }\right\Vert =1}\boldsymbol{%
\beta }^{\T }\boldsymbol{Q}_{\boldsymbol{\perp }}^{\T }\boldsymbol{D}%
\left( \xi _{n}\right) \boldsymbol{Q}\left( \boldsymbol{Q}^{\T }%
\boldsymbol{D}\left( \xi _{n}\right) \boldsymbol{Q}\right) ^{-2}\boldsymbol{Q%
}^{\T }\boldsymbol{D}\left( \xi _{n}\right) \boldsymbol{Q}_{\boldsymbol{%
\perp }}{\boldsymbol{\beta }}\right) .  \label{I2}
\end{equation}%
Noting that the maximum eigenvalue of a matrix $\boldsymbol{PP}^{\T }$
is that of $\boldsymbol{P}^{\T }\boldsymbol{P}$, and that $\boldsymbol{Q}%
_{\boldsymbol{\perp }}\boldsymbol{Q}_{\boldsymbol{\perp }}^{\T }=\left( 
\boldsymbol{I}_{N}-\boldsymbol{QQ}^{\T }\right) $, we have 
\begin{eqnarray*}
&&\max_{\left\Vert \boldsymbol{\beta }\right\Vert =1}\boldsymbol{\beta }%
^{\T }\boldsymbol{Q}_{\boldsymbol{\perp }}^{\T }\boldsymbol{D}\left(
\xi _{n}\right) \boldsymbol{Q}\left( \boldsymbol{Q}^{\T }\boldsymbol{D}%
\left( \xi _{n}\right) \boldsymbol{Q}\right) ^{-2}\boldsymbol{Q}^{\T }%
\boldsymbol{D}\left( \xi _{n}\right) \boldsymbol{Q}_{\boldsymbol{\perp }}%
\boldsymbol{\beta } \\
&=&ch_{\max }\boldsymbol{Q}_{\boldsymbol{\perp }}^{\T }\boldsymbol{D}%
\left( \xi _{n}\right) \boldsymbol{Q}\left( \boldsymbol{Q}^{\T }%
\boldsymbol{D}\left( \xi _{n}\right) \boldsymbol{Q}\right) ^{-1}\cdot \left( 
\boldsymbol{Q}^{\T }\boldsymbol{D}\left( \xi _{n}\right) \boldsymbol{Q}%
\right) ^{-1}\boldsymbol{Q}^{\T }\boldsymbol{D}\left( \xi _{n}\right) 
\boldsymbol{Q}_{\boldsymbol{\perp }} \\
&=&ch_{\max }\left( \boldsymbol{Q}^{\T }\boldsymbol{D}\left( \xi
_{n}\right) \boldsymbol{Q}\right) ^{-1}\boldsymbol{Q}^{\T }\boldsymbol{D}%
\left( \xi _{n}\right) \boldsymbol{Q}_{\boldsymbol{\perp }}\boldsymbol{Q}_{%
\boldsymbol{\perp }}^{\T }\boldsymbol{D}\left( \xi _{n}\right) 
\boldsymbol{Q}\left( \boldsymbol{Q}^{\T }\boldsymbol{D}\left( \xi
_{n}\right) \boldsymbol{Q}\right) ^{-1} \\
&=&ch_{\max }\left( \boldsymbol{Q}^{\T }\boldsymbol{D}\left( \xi
_{n}\right) \boldsymbol{Q}\right) ^{-1}\boldsymbol{Q}^{\T }\boldsymbol{D}%
^{2}\left( \xi _{n}\right) \boldsymbol{Q%
}\left( \boldsymbol{Q}^{\T }\boldsymbol{D}\left( \xi _{n}\right) 
\boldsymbol{Q}\right) ^{-1}-1.
\end{eqnarray*}%
This in (\ref{I2}) gives (\ref{I3}).\hfill\ $\square $

\subsection{\protect\bigskip Proof of Lemma \protect\ref{lemma: maxdiff}}

We calculate that, using the Huber score function, 
\begin{equation*}
\frac{\sigma _{M}^{2}}{\sigma ^{2}}=\frac{1-2c\phi \left( c\right) +2\left(
c^{2}-1\right) \Phi \left( -c\right) }{\left( 1-2\Phi \left( -c\right)
\right) ^{2}}\overset{def}{=}G\left( c\right) ,
\end{equation*}%
in terms of which $\nu _{\text{\textsc{m}}}=\left( \gamma ^{2}G\left(
c\right) +1\right) ^{-1}$. The function $G\ $is the asymptotic variance of
the M-estimate of location with $N(0,1)$ errors. It is decreasing in $c$ and
must exceed the asymptotic variance of the sample mean: $G\left( c\right)
>G\left( \infty \right) =1$. Two applications of L'H\^{o}pital's rule give $%
G\left( 0\right) =\pi /2$, the asymptotic variance of the median of a normal
sample, with corresponding regression estimate given by $\psi \left(
x\right) =sign(x)$. Thus 
\begin{equation*}
0<\nu _{\text{\textsc{ls}}}-\nu _{\text{\textsc{m}}}=\frac{\gamma ^{2}}{%
\gamma ^{2}+1}\left\{ \frac{G\left( c\right) -1}{\gamma ^{2}G\left( c\right)
+1}\right\} \overset{def}{=}H\left( c,\gamma ^{2}\right) .
\end{equation*}%
For fixed $\gamma $, $H\left( c,\gamma ^{2}\right) $ is increasing in $%
G\left( c\right) $, so is maximized at $c=0$. Then $H\left( 0,\gamma
^{2}\right) $ vanishes at $\gamma ^{2}=0,\infty $ and is maximized at $%
\gamma ^{2}=\left[ G\left( 0\right) \right] ^{-1/2}$, with value $%
\max_{c,\gamma }\left( \nu _{\text{\textsc{ls}}}-\nu _{\text{\textsc{m}}%
}\right) =\left( \sqrt{G\left( 0\right) }-1\right) \left/ \left( \sqrt{%
G\left( 0\right) }+1\right) \right. $. At the maximum, $\nu _{\text{\textsc{%
ls}}}=1/\left( \gamma ^{2}+1\right) $ and $\nu _{\text{\textsc{m}}}=1/\left(
\gamma ^{2}G\left( 0\right) +1\right) $, giving (\ref{nus}) . \hfill\ $%
\square $

\section*{Acknowledgements}

This work was carried out with the support of the Natural Sciences and
Engineering Council of Canada. It has benefited from the insightful comments of two anonymous referees.

\bibliographystyle{natbib}
\bibliography{references}


\end{document}
